\section{Method}\label{sec:method}

\paragraph{Roadmap.}
We begin with a length-only model that recovers $\Delta L$ weighting method; we then present a heteroscedastic Gaussian model with flexible priors, derive the posterior-precision aggregator, identify Jeffreys’ baseline ($\alpha{=}1$), and show how correctness-linked variance reduction yields $\alpha{<}1$; dependence corrections are incorporated via $s(L;\xi)$.

\subsection{Length-Only Benchmark and the \texorpdfstring{$\Delta L$}{Delta L}  Family}\label{subsec:warmup}
Under an i.i.d.\ token model with $\mathrm{Var}(\boldsymbol{Z}_{i,t})=\boldsymbol{\Sigma}_0$ and $\boldsymbol{\Gamma}_i(k)=\boldsymbol{0}$ for $k\neq 0$,
\begin{equation}
  \mathrm{Var}(\boldsymbol{g}_i)=L_i\,\boldsymbol{\Sigma}_0
  \quad\Longrightarrow\quad
  v_i \propto L_i.
\end{equation}
Substituting into \eqref{eq:mvu-weights} yields the inverse-length rule
\begin{equation}
  x_i \ \propto\ L_i^{-1},
  \label{eq:deltaL_alpha1}
\end{equation}
which is precisely the $\Delta L$ normalization of \citet{he2025deltal} with exponent $\alpha=1$. More generally, following \citet{he2025deltal} we consider the one-parameter family
\begin{equation}
  x_i \ \propto\ L_i^{-\beta},\qquad \beta\in[0,2],
  \label{eq:ldependence}
\end{equation}
interpolating between constant weighting ($\beta=0$) and inverse-length weighting ($\beta=1$), as a practical control of variance induced by heterogeneous lengths.

\paragraph{A two-parameter normalization family.}
Length-only normalization can be viewed as choosing a (possibly biased) \emph{global} step scale and a (possibly variance-reducing) \emph{relative} within-batch reweighting. To decouple these effects, we introduce a two-parameter family
\begin{equation}
  x_i(\alpha,\kappa)
  \;:=\;
  \frac{1}{M_\kappa}\,\frac{L_i^{-\alpha}}{\sum_{j=1}^{G} L_j^{-\alpha}},
  \qquad
  M_\kappa\;:=\;\Big(\sum_{j=1}^{G} L_j^{-\alpha}\Big)^{\kappa},
  \label{eq:two_param_norm}
\end{equation}
where $\alpha\in[0,1]$ controls the relative suppression of long responses and $\kappa\in\mathbb{R}$ controls the overall step magnitude. The standard $\Delta L$ estimator corresponds to $\kappa=0$ (purely relative reweighting with a fixed $M$), while $\kappa=1$ recovers a DAPO-style batch-length normalization where the step magnitude shrinks with the batch effective length. Section~\ref{sec:eb-theory} shows that CAPO preserves unbiasedness by enforcing the linear constraint $\sum_i x_i=1/M$ while learning the variance structure; empirically, we therefore treat $(\alpha,\kappa)$ as a diagnostic lens rather than as free hyperparameters for CAPO.

\subsection{Bayesian Setup under a General Prior}\label{subsec:bayes-setup}
We project each trajectory-level gradient onto a fixed direction to obtain a scalar summary. Let $\bm{u}$ be a unit vector and define
\begin{equation}
  g_i \;=\; \bm{u}^\top \bm{g}_i,\qquad \mathbb{E}[g_i]=\mu,\qquad \mathrm{Var}(g_i)=v_i.
  \label{eq:def_gi_scalar}
\end{equation}
\paragraph{Vector-valued gradients and common-shape covariances.}
Our Bayesian and empirical-Bayes derivations proceed via a scalar projection \eqref{eq:def_gi_scalar}.
This is not merely a notational convenience: in a broad and practically relevant regime, the resulting weights are also optimal for \emph{vector-valued} aggregation objectives.

\begin{assumption}[Common-shape covariance]\label{assump:common_shape}
There exists a fixed positive-semidefinite matrix $\boldsymbol{\Sigma}_0\in\mathbb{S}_{+}^{p}$ and positive scalars $v_i$ such that, for every trajectory $i$,
\[
\mathrm{Cov}(\boldsymbol{g}_i)=v_i\,\boldsymbol{\Sigma}_0.
\]
\end{assumption}

\begin{proposition}[Inverse-variance weights are optimal under common shape]\label{prop:common_shape_weights}
Assume cross-trajectory uncorrelatedness (Assumption~\ref{assump:cross_traj_uncorr}) and the common-shape model in Assumption~\ref{assump:common_shape}.
Let $\Psi:\mathbb{S}_{+}^{p}\to(0,\infty)$ be any positive-homogeneous functional, i.e.,
$\Psi(c\boldsymbol{A})=c\,\Psi(\boldsymbol{A})$ for all $c\ge 0$ and $\boldsymbol{A}\in\mathbb{S}_{+}^{p}$.
Then the unique minimizer of $\Psi(\mathrm{Cov}(\widehat{\boldsymbol{g}}))$ over weights satisfying $\sum_{i=1}^G x_i=1/M$ is given by the inverse-variance rule
$x_i^\star\propto v_i^{-1}$ (with normalization fixed by $\sum_i x_i=1/M$), i.e., the same weights as in Lemma~\ref{lem:mvu}.
\end{proposition}

\begin{proof}
Under cross-trajectory uncorrelatedness,
\[
\mathrm{Cov}(\widehat{\boldsymbol{g}})
=\sum_{i=1}^G x_i^2\,\mathrm{Cov}(\boldsymbol{g}_i)
=\sum_{i=1}^G x_i^2\,v_i\,\boldsymbol{\Sigma}_0
=\left(\sum_{i=1}^G x_i^2 v_i\right)\boldsymbol{\Sigma}_0.
\]
By positive homogeneity of $\Psi$,
\[
\Psi\!\left(\mathrm{Cov}(\widehat{\boldsymbol{g}})\right)
=\left(\sum_{i=1}^G x_i^2 v_i\right)\Psi(\boldsymbol{\Sigma}_0).
\]
Since $\Psi(\boldsymbol{\Sigma}_0)>0$ is constant in $(x_i)$, minimizing $\Psi(\mathrm{Cov}(\widehat{\boldsymbol{g}}))$ is equivalent to minimizing
$\sum_i x_i^2 v_i$ under $\sum_i x_i=1/M$, whose unique minimizer is the inverse-variance rule in Lemma~\ref{lem:mvu}.
\end{proof}

\begin{remark}[Beyond common shape]
When the covariance matrices $\mathrm{Cov}(\boldsymbol{g}_i)$ are not approximately proportional to a common shape,
the optimal weights can depend on the chosen scalarization objective $\Psi$ (e.g., a particular direction $\bm{u}$ or a trace objective).
CAPO focuses on the setting where a \emph{single} scalar proxy $v_i$ provides a stable and implementable approximation.
\end{remark}

Conditional on the variance parameters $\vartheta := (\sigma^2, \beta, \xi)$,
\begin{equation}
  g_i \,\big|\, \mu,\vartheta \ \sim\ \mathcal{N}\!\big(\mu,\ v_i(\vartheta)\big),
  \qquad
  v_i(\vartheta)\;=\;\sigma^2\,L_i^{\beta}\,s(L_i;\,\xi),
  \label{eq:hetero_general_prior}
\end{equation}
where $\sigma^2>0$ is a global scale, $\beta\in[0,2]$ controls length scaling, and $s(L_i;\xi)>0$ is a shape factor encoding within-trajectory dependence. This modular treatment allows us to plug in different estimates for the shape factor neatly.

\paragraph{General priors.}
We adopt a flat prior on the mean and a flexible, proper prior on the variance and shape parameters:
\begin{equation}
  p(\mu)\ \propto\ 1,\qquad
  \sigma^2 \sim \mathrm{Inv\text{-}Gamma}(a_0,b_0),\qquad
  \beta \sim \pi_\beta,\qquad
  \xi \sim \pi_\xi.
  \label{eq:priors_general}
\end{equation}
Jeffreys' prior for scale is recovered in the limit $a_0,b_0\to 0$. The prior $\pi_\beta$ can be chosen agnostic (centered at $1$ with moderate variance) or informative to reflect anticipated sub/superlinear growth. The prior $\pi_\xi$ encodes beliefs about token dependence.

\paragraph{Likelihood geometry and sufficient sums.}
Define precisions and aggregated quantities
\begin{equation}
  \lambda_i(\vartheta) := \frac{1}{v_i(\vartheta)} = \frac{1}{\sigma^2}\,\omega_i(\beta,\xi),
  \qquad
  \omega_i(\beta,\xi) := L_i^{-\beta}\,s(L_i;\xi)^{-1},
  \label{eq:def_lambda_omega}
\end{equation}
\begin{equation}
  \Lambda(\vartheta) \;=\; \sum_{i=1}^G \lambda_i(\vartheta),\qquad
  T(\vartheta) \;=\; \sum_{i=1}^G \lambda_i(\vartheta)\,g_i.
  \label{eq:Lambda_T}
\end{equation}
Then the conditional posterior for $\mu$ given $\vartheta$ is Normal:
\begin{equation}
  \mu\,\big|\,\vartheta, g_i \ \sim\ \mathcal{N}\!\big(m(\vartheta),\,\Lambda(\vartheta)^{-1}\big),
  \qquad
  m(\vartheta)=\frac{T(\vartheta)}{\Lambda(\vartheta)}=\frac{\sum_i \lambda_i(\vartheta) g_i}{\sum_i \lambda_i(\vartheta)}.
  \label{eq:mu_conditional}
\end{equation}
Let $m(\beta,\xi)$ denote $m(\vartheta)$ with $\sigma^2$ factored out via \eqref{eq:def_lambda_omega}, and define the weighted residual sum of squares
\begin{equation}
  \mathrm{RSS}_\omega(\beta,\xi) \;=\; \sum_{i=1}^G \omega_i(\beta,\xi)\,\big(g_i - m(\beta,\xi)\big)^2.
  \label{eq:rss_omega}
\end{equation}
Combining \eqref{eq:priors_general}--\eqref{eq:rss_omega}, the conditional posterior of $\sigma^2$ is
\begin{equation}
  \sigma^2 \,\big|\, \beta,\xi,\{g_i\} \ \sim\ \mathrm{Inv\text{-}Gamma}\!\left(a_0+\frac{G-1}{2},\ b_0+\frac{1}{2}\,\mathrm{RSS}_\omega(\beta,\xi)\right).
  \label{eq:sigma2_posterior}
\end{equation}

\begin{proposition}[Conditional posterior of $\sigma^2$]\label{prop:sigma2-posterior}
Under the Gaussian model
\[
g_i \mid \mu,\vartheta \ \sim\ \mathcal{N}\!\big(\mu,\ v_i(\vartheta)\big),
\qquad
v_i(\vartheta)=\sigma^2 L_i^{\beta}s(L_i;\xi)\ =\ \frac{\sigma^2}{\omega_i(\beta,\xi)},
\]
with $\vartheta=(\sigma^2,\beta,\xi)$ and $\omega_i(\beta,\xi):=L_i^{-\beta}s(L_i;\xi)^{-1}$ (cf.\ \eqref{eq:def_lambda_omega}),
and priors $p(\mu)\propto 1$, $\ \sigma^2\sim\mathrm{Inv\text{-}Gamma}(a_0,b_0)$ as in \eqref{eq:priors_general},
the conditional posterior of $\sigma^2$ given $(\beta,\xi)$ and data $\{g_i\}_{i=1}^G$ is
\begin{equation}
  \sigma^2 \,\big|\, \beta,\xi,\{g_i\}
  \ \sim\
  \mathrm{Inv\text{-}Gamma}\!\left(
    a_0+\frac{G-1}{2},\
    b_0+\frac{1}{2}\,\mathrm{RSS}_\omega(\beta,\xi)
  \right),
  \tag{\ref{eq:sigma2_posterior}}
\end{equation}
where the weighted mean is
\[
m(\beta,\xi)\ :=\ \frac{\sum_{i=1}^G \omega_i(\beta,\xi)\,g_i}{\sum_{i=1}^G \omega_i(\beta,\xi)}
\]
and the weighted residual sum of squares is
\[
\mathrm{RSS}_\omega(\beta,\xi)\ :=\ \sum_{i=1}^G \omega_i(\beta,\xi)\,\big(g_i-m(\beta,\xi)\big)^2.
\tag{\ref{eq:rss_omega}}
\]
\end{proposition}

\begin{proof}
\emph{Step 1: Likelihood in $(\mu,\sigma^2)$ up to $(\beta,\xi)$-constants.}
Conditional on $(\beta,\xi)$, the variances are $v_i=\sigma^2/\omega_i$.
Thus the likelihood is
\[
p(\{g_i\}\mid \mu,\sigma^2,\beta,\xi)
=\prod_{i=1}^G \frac{\sqrt{\omega_i}}{\sqrt{2\pi\sigma^2}}
\exp\!\left\{-\frac{\omega_i}{2\sigma^2}(g_i-\mu)^2\right\}.
\]
Dropping factors independent of $(\mu,\sigma^2)$ (i.e., $\prod_i \sqrt{\omega_i}$), the kernel is
\begin{equation}
p(\{g_i\}\mid \mu,\sigma^2,\beta,\xi)\ \propto\
(\sigma^2)^{-G/2}\,
\exp\!\left\{-\frac{1}{2\sigma^2}\sum_{i=1}^G \omega_i(g_i-\mu)^2\right\}.
\label{eq:lik-kernel}
\end{equation}

\emph{Step 2: Multiply by priors.}
With $p(\mu)\propto 1$ and $\sigma^2\sim\mathrm{Inv\text{-}Gamma}(a_0,b_0)$ whose density kernel is
$(\sigma^2)^{-a_0-1}\exp\{-b_0/\sigma^2\}$, the joint posterior kernel is
\[
p(\mu,\sigma^2\mid \beta,\xi,\{g_i\})
\ \propto\
(\sigma^2)^{-\left(a_0+1+\frac{G}{2}\right)}
\exp\!\left\{-\frac{1}{\sigma^2}\left(
b_0+\frac{1}{2}\sum_{i=1}^G \omega_i(g_i-\mu)^2
\right)\right\}.
\]

\emph{Step 3: Complete the square in $\mu$.}
Let $\Lambda_\omega:=\sum_{i=1}^G \omega_i$ and
$m(\beta,\xi):=\big(\sum_i \omega_i g_i\big)/\Lambda_\omega$.
Using the weighted-mean identity,
\[
\sum_{i=1}^G \omega_i(g_i-\mu)^2
=\sum_{i=1}^G \omega_i\big(g_i-m(\beta,\xi)\big)^2
+\Lambda_\omega\,(\mu-m(\beta,\xi))^2,
\]
we obtain
\[
p(\mu,\sigma^2\mid\beta,\xi,\{g_i\})
\ \propto\
(\sigma^2)^{-\left(a_0+1+\frac{G}{2}\right)}
\exp\!\left\{-\frac{1}{\sigma^2}\left(
b_0+\frac{1}{2}\mathrm{RSS}_\omega(\beta,\xi)
\right)\right\}\,
\exp\!\left\{-\frac{\Lambda_\omega}{2\sigma^2}\,(\mu-m(\beta,\xi))^2\right\}.
\]

\emph{Step 4: Integrate out $\mu$.}
Because $p(\mu)\propto 1$ and the $\mu$-kernel is Gaussian, we have
\[
\int_{\mathbb{R}}
\exp\!\left\{-\frac{\Lambda_\omega}{2\sigma^2}(\mu-m)^2\right\}\,d\mu
=\sqrt{\frac{2\pi\sigma^2}{\Lambda_\omega}},
\]
which contributes a factor $(\sigma^2)^{1/2}$ (and a constant in $(\beta,\xi)$) to the kernel.
Therefore,
\[
p(\sigma^2\mid \beta,\xi,\{g_i\})
\ \propto\
(\sigma^2)^{-\left(a_0+1+\frac{G}{2}\right)}(\sigma^2)^{1/2}\,
\exp\!\left\{-\frac{1}{\sigma^2}\left(
b_0+\frac{1}{2}\mathrm{RSS}_\omega(\beta,\xi)
\right)\right\}.
\]

\emph{Step 5: Identify the inverse-gamma form.}
Collecting powers of $\sigma^2$ gives
\[
p(\sigma^2\mid \beta,\xi,\{g_i\})
\ \propto\
(\sigma^2)^{-\left(a_0+\frac{G-1}{2}+1\right)}
\exp\!\left\{-\frac{1}{\sigma^2}\left(
b_0+\frac{1}{2}\mathrm{RSS}_\omega(\beta,\xi)
\right)\right\},
\]
which is the kernel of
\(
\mathrm{Inv\text{-}Gamma}\!\big(a_0+\tfrac{G-1}{2},\ b_0+\tfrac{1}{2}\mathrm{RSS}_\omega(\beta,\xi)\big).
\)
\end{proof}

\begin{remark}
(i) The factorization $\lambda_i(\vartheta)=\sigma^{-2}\omega_i(\beta,\xi)$ implies $m(\beta,\xi)$ does \emph{not} depend on $\sigma^2$, since the common factor cancels: $m(\vartheta)=T(\vartheta)/\Lambda(\vartheta)=\big(\sum \lambda_i g_i\big)/\big(\sum \lambda_i\big)=\big(\sum \omega_i g_i\big)/\big(\sum \omega_i\big)$. \\
(ii) The shape update $a_0\mapsto a_0+\frac{G-1}{2}$ reflects the loss of one degree of freedom from estimating the (flat-prior) mean $\mu$. \\
(iii) Constants that depend only on $(\beta,\xi)$, such as $\prod_i \sqrt{\omega_i}$ and $\Lambda_\omega^{-1/2}$, are immaterial for $p(\sigma^2\mid\beta,\xi,\{g_i\})$ and are absorbed into the normalizing constant.
\end{remark}

\subsection{Formal Results and Aggregators}\label{subsec:formal-results}

\begin{lemma}[MVU weights with known $v_i$]\label{lem:mvu}
Let $(g_1,\dots,g_G)$ be scalar random variables such that $\mathbb{E}[g_i]=\mu$ for all $i$,
$\mathrm{Var}(g_i)=v_i\in(0,\infty)$ are known, and the trajectories are cross-uncorrelated:
$\mathrm{Cov}(g_i,g_j)=0$ for all $i\neq j$.
For any weights $(x_1,\dots,x_G)$ satisfying the unbiasedness constraint $\sum_{i=1}^G x_i = 1/M$,
the linear estimator $\widehat{\mu}:=\sum_{i=1}^G x_i g_i$ is unbiased for $\mu/M$ and satisfies
\[
\mathrm{Var}(\widehat{\mu})=\sum_{i=1}^G x_i^2 v_i.
\]
The unique minimizer of $\mathrm{Var}(\widehat{\mu})$ over all such unbiased linear estimators is
\begin{equation}
  x_i^\star \;=\; \frac{1}{M}\,\frac{v_i^{-1}}{\sum_{j=1}^G v_j^{-1}},
  \tag{\ref{eq:mvu-weights} revisited}
\end{equation}
i.e., inverse-variance weighting with normalization fixed by $\sum_i x_i=1/M$.
\end{lemma}

\begin{proof}
Under the stated cross-uncorrelatedness, for any weights $(x_1,\dots,x_G)$ we have
\[
\mathrm{Var}\!\left(\sum_{i=1}^G x_i g_i\right)
=\sum_{i=1}^G x_i^2\,\mathrm{Var}(g_i)
=\sum_{i=1}^G x_i^2 v_i.
\]
We minimize the strictly convex quadratic function $Q(x):=\sum_{i=1}^G x_i^2 v_i$
over the affine subspace $\{x\in\mathbb{R}^G:\sum_{i=1}^G x_i=1/M\}$.

Introduce the Lagrangian
\[
\mathcal{L}(x_1,\dots,x_G,\lambda)
=\sum_{i=1}^G x_i^2 v_i
-\lambda\left(\sum_{i=1}^G x_i-\frac{1}{M}\right).
\]
The first-order optimality conditions are, for each $i\in\{1,\dots,G\}$,
\[
\frac{\partial \mathcal{L}}{\partial x_i}=2x_i v_i-\lambda=0
\quad\Longrightarrow\quad
x_i=\frac{\lambda}{2v_i}.
\]
Imposing the constraint $\sum_i x_i=1/M$ yields
\[
\frac{1}{M}=\sum_{i=1}^G x_i=\frac{\lambda}{2}\sum_{i=1}^G v_i^{-1}
\quad\Longrightarrow\quad
\frac{\lambda}{2}=\frac{1/M}{\sum_{j=1}^G v_j^{-1}}.
\]
Substituting back gives
\[
x_i^\star=\frac{1}{M}\,\frac{v_i^{-1}}{\sum_{j=1}^G v_j^{-1}},
\]
which satisfies the constraint and the first-order conditions.
Since $Q$ is strictly convex (all $v_i>0$), the solution is the unique global minimizer.
\end{proof}

\begin{proposition}[Robustness of plug-in inverse-variance weights]\label{prop:plugin_robust}
Assume the setting of Lemma~\ref{lem:mvu} and let $\{\widehat{v}_i\}_{i=1}^G$ be positive variance estimates such that, for some $c\ge 1$,
\begin{equation}
  \frac{1}{c}\,v_i \;\le\; \widehat{v}_i \;\le\; c\,v_i
  \qquad \text{for all } i\in\{1,\dots,G\}.
  \label{eq:vi-mult-error}
\end{equation}
Define the plug-in inverse-variance weights
\[
\widehat{x}_i \;:=\; \frac{1}{M}\,\frac{\widehat{v}_i^{-1}}{\sum_{j=1}^G \widehat{v}_j^{-1}}.
\]
Let $V^\star$ denote the minimum variance achieved by the oracle MVU weights of Lemma~\ref{lem:mvu}, and let
$\widehat{V}:=\mathrm{Var}(\sum_i \widehat{x}_i g_i)$ denote the variance under the plug-in weights.
Then the plug-in rule is variance-robust in the sense that
\[
\frac{1}{c^2}\,V^\star \;\le\; \widehat{V} \;\le\; c^2\,V^\star.
\]
\end{proposition}

\begin{proof}
Let $a_i:=v_i^{-1}$ and $\widehat{a}_i:=\widehat{v}_i^{-1}$ so that \eqref{eq:vi-mult-error} is equivalent to
$\widehat{a}_i\in[a_i/c,\,c a_i]$.
Define $S:=\sum_{i=1}^G a_i$ and $\widehat{S}:=\sum_{i=1}^G \widehat{a}_i$.
A direct computation (see the proof of Lemma~\ref{lem:mvu}) gives the oracle optimum
\[
V^\star=\frac{1}{M^2}\,\frac{1}{S}.
\]
For the plug-in weights, cross-trajectory uncorrelatedness gives
\[
\widehat{V}
=\sum_{i=1}^G \widehat{x}_i^2 v_i
=\frac{1}{M^2}\,\frac{\sum_{i=1}^G \widehat{a}_i^2 v_i}{\widehat{S}^2}
=\frac{1}{M^2}\,\frac{\sum_{i=1}^G \widehat{a}_i^2/a_i}{\widehat{S}^2}.
\]
Since $\widehat{a}_i/a_i\le c$, we have $\widehat{a}_i^2/a_i=\widehat{a}_i(\widehat{a}_i/a_i)\le c\,\widehat{a}_i$ and hence
$\sum_i \widehat{a}_i^2/a_i \le c\,\widehat{S}$.
Therefore
\[
\widehat{V}\le \frac{1}{M^2}\,\frac{c}{\widehat{S}}.
\]
Moreover, $\widehat{a}_i\ge a_i/c$ implies $\widehat{S}\ge S/c$, so $\widehat{V}\le \frac{1}{M^2}\frac{c}{S/c}=c^2 V^\star$.

For the lower bound, $\widehat{a}_i/a_i\ge 1/c$ implies $\sum_i \widehat{a}_i^2/a_i \ge (1/c)\widehat{S}$, hence
$\widehat{V}\ge \frac{1}{M^2}\frac{1/c}{\widehat{S}}$.
Since $\widehat{a}_i\le c a_i$ implies $\widehat{S}\le c S$, we conclude
$\widehat{V}\ge \frac{1}{M^2}\frac{1/c}{c S}=\frac{1}{c^2}V^\star$.
\end{proof}


\begin{theorem}[Posterior-precision aggregation and scale cancellation]\label{thm:ppa}
Consider the scalarized Gaussian model \eqref{eq:def_gi_scalar}--\eqref{eq:hetero_general_prior} and assume
the heteroscedastic variances factorize as
\begin{equation}
  v_i(\vartheta)=\sigma^2\,f_i(\vartheta),
  \qquad f_i(\vartheta)>0,
  \label{eq:vi-factor}
\end{equation}
so that the conditional precisions satisfy $\lambda_i(\vartheta)=1/v_i(\vartheta)=(1/\sigma^2)\,f_i(\vartheta)^{-1}$.
Fix $\vartheta$ (e.g., $\vartheta=(\beta,\xi)$) and endow $\sigma^2$ with any prior such that
$\mathbb{E}[\sigma^{-2}\mid \vartheta,\{g_j\}_{j=1}^G]$ exists and is finite.

Define the \emph{posterior-precision scores}
\begin{equation}
  \widetilde{w}_i(\vartheta)
  \;:=\;
  \mathbb{E}\!\left[\lambda_i(\vartheta)\,\middle|\,\vartheta,\{g_j\}_{j=1}^G\right]
  \;=\;
  f_i(\vartheta)^{-1}\,\mathbb{E}\!\left[\sigma^{-2}\,\middle|\,\vartheta,\{g_j\}_{j=1}^G\right].
  \label{eq:posterior-precision-scores}
\end{equation}
Then the (unique) minimum-variance unbiased linear estimator of $\mu/M$ of the form
$\widehat{\mu}=\sum_{i=1}^G x_i g_i$ under the constraint $\sum_{i=1}^G x_i=1/M$ is obtained by
normalizing the relative precisions:
\begin{equation}
  x_i^{\mathrm{Bayes}}(\vartheta)
  \;=\;
  \frac{1}{M}\,\frac{f_i(\vartheta)^{-1}}{\sum_{j=1}^G f_j(\vartheta)^{-1}}
  \;=\;
  \frac{1}{M}\,\frac{\widetilde{w}_i(\vartheta)}{\sum_{j=1}^G \widetilde{w}_j(\vartheta)}.
  \label{eq:bayes_weight_def}
\end{equation}
In particular, the common scale factor $\mathbb{E}[\sigma^{-2}\mid\vartheta,\{g_j\}]$ cancels from the normalized weights.
\end{theorem}

\begin{proof}
Fix $\vartheta$ and condition on $(\sigma^2,\vartheta)$.
Under \eqref{eq:vi-factor}, the conditional variances are $v_i(\vartheta)=\sigma^2 f_i(\vartheta)$ and hence
\[
v_i(\vartheta)^{-1}=\lambda_i(\vartheta)=\frac{1}{\sigma^2}\,f_i(\vartheta)^{-1}.
\]
Therefore, conditional on $(\sigma^2,\vartheta)$ the variances $\{v_i(\vartheta)\}$ are known and
Lemma~\ref{lem:mvu} applies, yielding the unique minimizer
\[
x_i^\star(\sigma^2,\vartheta)
=\frac{1}{M}\,\frac{v_i(\vartheta)^{-1}}{\sum_{j=1}^G v_j(\vartheta)^{-1}}
=\frac{1}{M}\,\frac{f_i(\vartheta)^{-1}}{\sum_{j=1}^G f_j(\vartheta)^{-1}},
\]
where the factor $\sigma^{-2}$ cancels from numerator and denominator.
This is exactly the first equality in \eqref{eq:bayes_weight_def}.

For the second equality, use \eqref{eq:posterior-precision-scores} to write
$\widetilde w_i(\vartheta)=f_i(\vartheta)^{-1}\,\mathbb{E}[\sigma^{-2}\mid\vartheta,\{g_j\}]$;
normalizing $\widetilde w_i(\vartheta)$ across $i$ cancels the common factor
$\mathbb{E}[\sigma^{-2}\mid\vartheta,\{g_j\}]$, giving the same weights as above.
\end{proof}

\begin{remark}[When $\vartheta$ is random]
If $\vartheta$ itself is endowed with a non-degenerate prior (e.g., a prior on the length exponent $\beta$),
then two distinct Bayesian constructions become relevant:
\begin{enumerate}
\item The \emph{exact} posterior mean of the conditional MVU weights
$x_i^\star(\vartheta)\propto f_i(\vartheta)^{-1}$, namely
$\mathbb{E}[x_i^\star(\vartheta)\mid\{g_j\}]$.
\item The \emph{posterior-precision approximation} that replaces $f_i(\vartheta)^{-1}$ by its posterior mean
$\mathbb{E}[f_i(\vartheta)^{-1}\mid\{g_j\}]$ and then renormalizes across $i$.
\end{enumerate}
These coincide whenever the posterior over $\vartheta$ is sufficiently concentrated or when $f_i(\vartheta)$
is treated as deterministic via Empirical Bayes (as done in \S\ref{sec:eb-theory}).
\end{remark}

\begin{theorem}[Jeffreys’ baseline recovers $\Delta L$ with $\alpha=1$]\label{thm:jeffreys}
Assume independence $s(\cdot;\xi)\equiv 1$, fix $\beta=1$, and take $a_0,b_0\to 0$ so that $p(\sigma^2)\propto 1/\sigma^2$. Then
\begin{equation}
  x_i^{\mathrm{Bayes}} \ \propto\ L_i^{-1}.
  \tag{\ref{eq:deltaL_alpha1} revisited}
\end{equation}
\end{theorem}

\begin{proof}
With $\beta=1$ and $s\equiv 1$, we have $\lambda_i(\vartheta)=\sigma^{-2} L_i^{-1}$. From \eqref{eq:sigma2_posterior}, $\mathbb{E}[\sigma^{-2}\mid \beta{=}1,\{g\}]=(G-1)/\mathrm{RSS}_\omega(\beta{=}1)$, a common factor that cancels in Theorem~\ref{thm:ppa}, yielding $x_i\propto L_i^{-1}$.
\end{proof}

\begin{proposition}[Correctness-linked contraction implies $\alpha<1$]\label{prop:alpha-less-than-1}
Suppose the per-token variance contracts with response length via $\kappa(L)\in(0,1]$ with $\kappa'(L)\le 0$ and
\(
\kappa(L)=c\,L^{-\delta}(1+o(1))
\)
for some $\delta\in[0,1)$, so that
\begin{equation}
  \mathrm{Var}(g_i\mid L_i=L)=\sigma^2\,L\,\kappa(L)\,s(L;\xi)
  \;=\;
  \sigma^2\,L^{1-\delta}\,s(L;\xi)\,(1+o(1)).
  \label{eq:effective-length-law_repeat}
\end{equation}
Then in the independence benchmark $s\equiv 1$ the implied $\Delta L$ exponent is
\(
\alpha=\beta=1-\delta<1.
\)
\end{proposition}

\begin{proof}
The effective length law in \eqref{eq:effective-length-law_repeat} gives variance exponent $\beta=1-\delta$. Applying Theorem~\ref{thm:ppa} with Jeffreys’ scale (or any scale prior whose common factor cancels) yields weights $\propto L^{-\beta}$, hence $\alpha=\beta=1-\delta<1$.
\end{proof}

\begin{corollary}[Serial dependence corrections]\label{cor:dependence}
If $s(L;\xi)=\Theta(1)$ (e.g., AR(1) or $k$-band dependence), then the large-$L$ exponent in Proposition~\ref{prop:alpha-less-than-1} is unchanged; dependence enters through the multiplicative factor $s(L;\xi)^{-1}$ in the weights.
\end{corollary}

\subsection{Informative Prior on the Length Exponent}\label{subsec:informative-beta}
Retain independence ($s\equiv 1$) but let $\beta\sim\pi_\beta$ with posterior $p(\beta\mid\{g_i\})$. Using \eqref{eq:bayes_weight_def},
\begin{equation}
  \widetilde{w}_i \;\propto\; \mathbb{E}\!\left[L_i^{-\beta}\,\big|\,\{g_j\}\right].
  \label{eq:wtilde_beta_random_repeat}
\end{equation}
Two useful consequences:
\begin{enumerate}
  \item \emph{MAP/Laplace approximation.} If $p(\beta\mid\{g_i\})$ is sharply peaked at $\hat\beta$, then
  \(
    x_i^{\mathrm{Bayes}} \propto L_i^{-\hat\beta}.
  \)
  \item \emph{Gaussian posterior moment.} If $\beta\mid\{g_i\}\sim \mathcal{N}(m_\beta,\tau_\beta^2)$ (ignoring truncation),
  \begin{equation}
    \mathbb{E}\!\left[L_i^{-\beta}\mid\{g_i\}\right]
    \;=\;
    \exp\!\Big(-m_\beta\ln L_i + \tfrac{1}{2}\tau_\beta^2(\ln L_i)^2\Big)
    \;\approx\; L_i^{-m_\beta}\quad (\tau_\beta^2\ \text{small}),
  \end{equation}
  i.e., an effective $\alpha\simeq m_\beta$.
\end{enumerate}

\subsection{Why \texorpdfstring{$\beta<1$}{beta < 1} When Longer Responses Are More Likely Correct}\label{subsec:alpha-less-than-1}
\paragraph{Intuition.}
If longer responses have a higher probability of being correct, then, after baselining, their token contributions are less noisy on average. In a heteroscedastic model, the per-token variance decreases with length, so the sum over $L$ tokens grows sub-linearly in $L$. Consequently, inverse-variance weights scale as $L^{-\alpha}$ with $\alpha<1$.

\paragraph{Heteroskedastic variance model tied to correctness.}
Let $g_i=\sum_{t=1}^{L_i} Z_{i,t}$ be the scalarized trajectory gradient (see \eqref{eq:traj-grad} and \eqref{eq:hetero_general_prior}). Suppose
\begin{equation}
  \mathrm{Var}(Z_{i,t}\mid L_i=L)\;=\;\sigma^2\,\kappa(L),
  \qquad
  \kappa'(L)\le 0,
\end{equation}
with $\kappa(L)\propto p(L)\{1-p(L)\}$ for a correctness probability $p(L)$ with $p(L)\uparrow 1$ as $L\uparrow\infty$. Then
\begin{equation}
  \mathrm{Var}(g_i\mid L_i=L)\;=\;\sigma^2\,L\,\kappa(L)\,s(L;\xi).
\end{equation}
If $\kappa(L)=c\,L^{-\delta}(1+o(1))$ with $\delta\in[0,1)$, we recover \eqref{eq:effective-length-law_repeat} and thus $\alpha=1-\delta<1$ by Proposition~\ref{prop:alpha-less-than-1}. The conclusion persists with serial dependence by Corollary~\ref{cor:dependence}.


\subsection{Weights under \texorpdfstring{$k$}{k}-banded Covariance Structure}\label{subsec:kband-eta}

We work with the scalarized series $Y_{i,t}:=\bm{u}^\top \boldsymbol{Z}_{i,t}$ and $g_i=\sum_{t=1}^{L_i}Y_{i,t}$ as in \eqref{eq:def_gi_scalar}. Assume second–order stationarity and the \emph{$k$-banded, stretched-geometric} auto-covariance:
\begin{equation}
\gamma(0)=\mathrm{Var}(Y_{i,t})=\tau^2,\qquad
\gamma(h)=\mathrm{Cov}(Y_{i,t},Y_{i,t-h})=\tau^2\,\rho^{\,|h|^{\eta}}\ \mathbf{1}\{|h|\le k\},
\quad \eta\ge 0,
\label{eq:kband-stretch}
\end{equation}
with $|\rho|<1$ for $\eta>0$. The case $\eta=1$ reduces to geometric decay $\gamma(h)=\tau^2\rho^{|h|}$; $\eta\downarrow 0$ yields \emph{no decay across lags} $\gamma(h)=\tau^2\rho$ (compound-symmetry within the band). Let $m_i:=\min\{k,L_i-1\}$.

\begin{lemma}[Dependence shape factor for stretched-geometric $k$-band]\label{lem:sL-kband-eta}
Under \eqref{eq:kband-stretch},
% \begin{equation}
% \mathrm{Var}(g_i)
% \;=\; \tau^2\!\left[
%   L_i \;+\; 2\sum_{h=1}^{m_i} (L_i-h)\,\rho^{\,h^{\eta}}
% \right]
% \;=\; \tau^2\,L_i\,s(L_i;\rho,k,\eta),
% \\
% s(L;\rho,k,\eta)\;=\;1+\frac{2}{L}\sum_{h=1}^{m} (L-h)\rho^{\,h^{\eta}}.
% \label{eq:v-and-s-kband-eta}
% \end{equation}

\begin{align}
\mathrm{Var}(g_i)
&= \tau^2\!\left[
  L_i + 2\sum_{h=1}^{m_i} (L_i-h)\,\rho^{\,h^{\eta}}
\right] \\
&= \tau^2\,L_i\,s(L_i;\rho,k,\eta),
\label{eq:v-and-s-kband-eta}
\end{align}

\begin{align}
s(L;\rho,k,\eta)
&= 1+\frac{2}{L}\sum_{h=1}^{m} (L-h)\rho^{\,h^{\eta}}.
\end{align}
\end{lemma}

\begin{proof}
By Bartlett’s identity (Lemma~\ref{lem:bartlett}),
$\mathrm{Var}(g_i)=\sum_{h=-(L_i-1)}^{L_i-1}(L_i-|h|)\gamma(h)$.
Insert \eqref{eq:kband-stretch}, collect positive and negative lags, and truncate at $m_i$.
\end{proof}

\begin{proposition}[Weighting rule under stretched-geometric $k$-band]\label{prop:weights-kband-eta}
Under the heteroscedastic model \eqref{eq:hetero_general_prior} with
\(
v_i(\vartheta)=\sigma^2 L_i^{\beta}\,s(L_i;\rho,k,\eta)
\),
the posterior-precision (Bayes) weights satisfy, up to a common factor,
\begin{equation}
x_i^{\mathrm{Bayes}} \ \propto\ \frac{1}{L_i^{\beta}\,s(L_i;\rho,k,\eta)}
\;=\;
\frac{1}{L_i^{\beta}\Big(1+\frac{2}{L_i}\sum_{h=1}^{m_i}(L_i-h)\rho^{\,h^{\eta}}\Big)}.
\label{eq:weights-kband-eta}
\end{equation}
In particular, under the Jeffreys baseline $\beta=1$,
\(
x_i \propto \big[L_i + 2\sum_{h=1}^{m_i}(L_i-h)\rho^{\,h^{\eta}}\big]^{-1}.
\)
\end{proposition}

\begin{proof}
From \eqref{eq:hetero_general_prior} and Lemma~\ref{lem:sL-kband-eta},
\(
\lambda_i(\vartheta)=1/v_i(\vartheta)=(1/\sigma^2)\,L_i^{-\beta}\,s(L_i;\rho,k,\eta)^{-1}
\).
Apply Theorem~\ref{thm:ppa} and note $1/\sigma^2$ cancels upon normalization.
\end{proof}

\begin{lemma}[Monotonicity in the decay exponent]\label{lem:mono-eta}
For fixed $(\rho,k,L)$ with $\rho\in(0,1)$, $s(L;\rho,k,\eta)$ is non-increasing in $\eta$:
\[
0\le \eta_1<\eta_2 \ \Rightarrow\ s(L;\rho,k,\eta_1)\ \ge\ s(L;\rho,k,\eta_2).
\]
Consequently, the weight magnitude $[L^{\beta}s(L;\rho,k,\eta)]^{-1}$ is non-decreasing in $\eta$.
\end{lemma}

\begin{proof}
For $h\ge 1$ and $\rho\in(0,1)$, the map $\eta\mapsto \rho^{\,h^{\eta}}$ is decreasing, hence each summand in \eqref{eq:v-and-s-kband-eta} decreases in $\eta$.
\end{proof}

\paragraph{No-decay limit ($\eta=0$): constant-lag covariance within the band.}
Setting $\eta=0$ in \eqref{eq:kband-stretch} gives $\gamma(h)=\tau^2\rho$ for $1\le|h|\le k$. Then
\begin{equation}
s(L;\rho,k,0)
\;=\;
1+\frac{2\rho}{L}\sum_{h=1}^{m}(L-h)
\;=\;
1+2m\rho-\frac{m(m+1)}{L}\rho,
\qquad m=\min\{k,L-1\}.
\label{eq:s-kband-nodecay}
\end{equation}
Accordingly,
\begin{equation}
x_i^{\mathrm{Bayes}}
\ \propto\
\frac{1}{L_i^{\beta}\big(1+2m_i\rho-\frac{m_i(m_i+1)}{L_i}\rho\big)},
\qquad m_i=\min\{k,L_i-1\}.
\label{eq:weights-kband-nodecay}
\end{equation}

\paragraph{Compound symmetry.}
Taking $k\ge L-1$ in \eqref{eq:s-kband-nodecay} yields
\(
m=L-1
\)
and
\begin{equation}
s(L;\rho,k\!\ge\!L-1,0)
=
1+\rho(L-1)
=
(1-\rho)+\rho L,
\label{eq:s-compound}
\end{equation}
which is the compound-symmetry shape from Lemma~\ref{lem:cs-var}. Thus
\begin{equation}
x_i^{\mathrm{Bayes}}
\ \propto\
\frac{1}{L_i^{\beta}\,\big(1+\rho(L_i-1)\big)}
\qquad\text{and, for }\beta=1,\qquad
x_i\ \propto\ \frac{1}{(1-\rho)L_i+\rho L_i^2}.
\label{eq:weights-compound}
\end{equation}

\begin{remark}[Asymptotics]
For fixed $k$ and $L\to\infty$,
\(
s(L;\rho,k,\eta)=1+2\sum_{h=1}^{k}\rho^{\,h^{\eta}}+O(L^{-1})
\),
a positive constant depending on $(\rho,k,\eta)$. Hence the length exponent is unchanged and dependence multiplies the $\Delta L$ rule by a constant. Under compound symmetry (no band), $s(L;\rho,\,\cdot,0)=(1-\rho)+\rho L$, so $v_i=\Theta(L_i^{\beta+1})$ and, under $\beta=1$, $x_i\sim ( \rho L_i^2 )^{-1}$ as $L_i\to\infty$.
\end{remark}


\subsection{Why \texorpdfstring{$\beta>1$}{beta > 1} When Longer Responses Are More Likely Correct (with tail degradation)}\label{subsec:alpha-greater-than-1}

\paragraph{Intuition.}
If longer responses are more likely to be correct up to a point, then after base-lining, their token contributions become less noisy as length grows toward an optimum length $L^\star$. Beyond $L^\star$, if the model gets lost, correctness drops and noise increases again. In a heteroskedastic model, this corresponds to a \emph{U-shaped} per-token variance profile $\kappa(L)$: initially decreasing, then increasing. Consequently, the variance of the sum over $L$ tokens grows \emph{sublinearly} before $L^\star$ and \emph{superlinearly} in the tail, yielding inverse-variance weights that scale like $L^{-\beta(L)}$ with $\beta(L)<1$ for moderate $L$ and $\beta(L)>1$ in the tail.

\paragraph{Correctness-weighted Heteroscedastic variance model (non-monotone variance profile).}
Let $g_i=\sum_{t=1}^{L_i} Z_{i,t}$ be the scalarized trajectory gradient (see \eqref{eq:traj-grad} and \eqref{eq:hetero_general_prior}). As before, write
\begin{equation}
  \mathrm{Var}(Z_{i,t}\mid L_i=L)\;=\;\sigma^2\,\kappa(L),
  \qquad
  \mathrm{Var}(g_i\mid L_i=L)\;=\;\sigma^2\,L\,\kappa(L)\,s(L;\xi),
  \label{eq:var-sum-kappa}
\end{equation}
where $s(L;\xi)$ encodes serial dependence (Sec.~\ref{subsec:kband-eta}). We now allow $\kappa(L)$ to be \emph{non-monotone}:
\begin{assumption}[Tail degradation]\label{assump:tail}
There exist $L^\star\in(1,\infty)$ and exponents $\delta\in[0,1)$, $\theta>0$, and constants $c_-,c_+>0$ such that
\[
\kappa(L)
=\begin{cases}
c_-\,L^{-\delta}\,\big(1+o(1)\big), & L\nearrow L^\star \text{ (moderate regime)},\\[2pt]
c_+\,L^{\ \theta}\,\big(1+o(1)\big), & L\to\infty \text{ (tail regime)}.
\end{cases}
\]
That is, $\kappa$ decreases polynomially up to $L^\star$ and increases polynomially beyond it.
\end{assumption}

\paragraph{Local length exponent and effective law.}
Define the \emph{local variance elasticity}
\[
\varepsilon(L)\ :=\ \frac{d\log \kappa(L)}{d\log L}.
\]
From \eqref{eq:var-sum-kappa},
\(
\log v(L)=\log\sigma^2+\log L+\log\kappa(L)+\log s(L;\xi)
\),
so (ignoring the slowly varying $s(L;\xi)$ for the exponent),
\begin{equation}
  \beta_{\mathrm{eff}}(L)
  \ :=\ \frac{d\log v(L)}{d\log L}\ \approx\ 1+\varepsilon(L).
  \label{eq:beta-local}
\end{equation}
Assumption~\ref{assump:tail} implies
\[
\beta_{\mathrm{eff}}(L)\approx
\begin{cases}
1-\delta, & L\lesssim L^\star,\\
1+\theta, & L\gg L^\star,
\end{cases}
\qquad
\text{so}\quad
v(L)\ \asymp\
\begin{cases}
\sigma^2\,L^{\,1-\delta}\,s(L;\xi), & L\lesssim L^\star,\\
\sigma^2\,L^{\,1+\theta}\,s(L;\xi), & L\gg L^\star.
\end{cases}
\label{eq:effective-length-piecewise}
\]

\begin{proposition}[Posterior-precision weights with tail degradation]\label{prop:tail-weights}
Under the heteroskedastic Normal model with variance $v(L)$ as in \eqref{eq:var-sum-kappa} and Assumption~\ref{assump:tail}, the Bayes (posterior-precision) weights satisfy, up to a common factor,
\begin{equation}
x_i^{\mathrm{Bayes}}
\ \propto\
\frac{1}{v(L_i)}
\ \approx\
\begin{cases}
L_i^{-(1-\delta)}\,s(L_i;\xi)^{-1}, & L_i\lesssim L^\star,\\[2pt]
L_i^{-(1+\theta)}\,s(L_i;\xi)^{-1}, & L_i\gg L^\star.
\end{cases}
\label{eq:weights-piecewise}
\end{equation}
Equivalently, the implied $\Delta L$ exponent is \(\beta(L)\approx\beta_{\mathrm{eff}}(L)\), i.e., $\beta(L)<1$ before $L^\star$ and $\beta(L)>1$ in the tail.
\end{proposition}

\begin{proof}
Theorem~\ref{thm:ppa} yields $x_i^{\mathrm{Bayes}}\propto\mathbb{E}[1/v(L_i)\mid\text{data}]$. Under Assumption~\ref{assump:tail}, $v(L)$ is regularly varying with index $1-\delta$ for $L\lesssim L^\star$ and $1+\theta$ for $L\to\infty$, hence the stated asymptotics. The multiplicative $s(L;\xi)$ enters as in Sec.~\ref{subsec:kband-eta}.
\end{proof}

\paragraph{Impact of serial dependence.}
The factor $s(L;\xi)$ multiplies $v(L)$ and hence divides the weights:
\[
x_i^{\mathrm{Bayes}}\ \propto\ \frac{1}{L_i\,\kappa(L_i)\,s(L_i;\xi)}.
\]
For $k$-banded dependence with fixed $k$ (Sec.~\ref{subsec:kband-eta}), $s(L;\xi)\to s_\infty(\xi)$, a constant; the exponents in \eqref{eq:effective-length-piecewise} remain $1\pm(\delta,\theta)$. Under compound symmetry (constant-lag covariance; Sec.~\ref{subsec:cs}), $s(L;\rho)=(1-\rho)+\rho L$ contributes an extra factor $\sim \rho L$ as $L\to\infty$, so the tail exponent becomes $2+\theta$ and
\[
x_i^{\mathrm{Bayes}}\ \propto\ L_i^{-(2+\theta)}\quad\text{(tail, compound symmetry)}.
\]

\paragraph{Concrete parametric example (unimodal correctness).}
Let the correctness probability be unimodal in $\log L$, e.g.,
\[
p(L)=\big(1+\exp\{-a_0-a_1\log L + a_2(\log L)^2\}\big)^{-1},
\qquad a_2>0,
\]
so $p(L)\uparrow$ for small $L$ and $p(L)\downarrow$ for large $L$. If the variance factor is $\kappa(L)\propto p(L)\{1-p(L)\}$ around the interior (maximal near $p=1/2$) but is \emph{augmented} by a tail noise term $\nu(L)\asymp L^{\theta}$ (credit-assignment drift), i.e.,
\[
\kappa(L)\ \asymp\ c_- L^{-\delta} \ \text{ for } L\lesssim L^\star,
\qquad
\kappa(L)\ \asymp\ c_+ L^{\theta} \ \text{ for } L\gg L^\star,
\]
then Proposition~\ref{prop:tail-weights} applies with the same exponents.

\begin{remark}
Practical summary.
    \begin{itemize}
    \item For \textbf{moderate lengths} where correctness improves, use \(\beta\simeq 1-\delta<1\).
    \item In the \textbf{long-tail} where the model gets lost, use \(\beta\simeq 1+\theta>1\) (and add $+1$ more under compound symmetry).
    \item The posterior-precision rule automatically enacts this via $x_i\propto [L_i\,\kappa(L_i)\,s(L_i;\xi)]^{-1}$; the common scale cancels.
\end{itemize}
\end{remark}

% =========================================================
% §3.X  Empirical Bayes estimator for length and dependence
% Place this AT THE END of your "Special Cases" subsection in Methods.
% Reference as \ref{sec:eb-theory}. The Algorithms go in a SEPARATE SECTION.
% =========================================================


\subsection{Empirical Bayes estimator for length and dependence}
\label{sec:eb-theory}

% --- Safe theorem/proof envs if not already defined in the preamble
\makeatletter
\@ifundefined{theorem}{\newtheorem{theorem}{Theorem}[section]}{}
\@ifundefined{lemma}{\newtheorem{lemma}[theorem]{Lemma}}{}
\@ifundefined{proposition}{\newtheorem{proposition}[theorem]{Proposition}}{}
\@ifundefined{corollary}{\newtheorem{corollary}[theorem]{Corollary}}{}
\@ifundefined{assumption}{\newtheorem{assumption}[theorem]{Assumption}}{}
\@ifundefined{definition}{\newtheorem{definition}[theorem]{Definition}}{}
\@ifundefined{remark}{\newtheorem{remark}[theorem]{Remark}}{}
\makeatother
\ifx\proof\undefined
  \newenvironment{proof}{\par\noindent\textit{Proof.}\ }{\hfill$\square$\par}
\fi

We consider trajectories \(i=1,\dots,G\) with length \(L_i\in\mathbb N\) and a scalarized trajectory-gradient statistic \(g_i\in\mathbb R\) (e.g., a fixed-direction projection of the per-trajectory policy-gradient, as in \eqref{eq:def_gi_scalar}).
Adopt the Normal working model
\begin{equation}
\label{eq:model}
g_i \mid \mu,\sigma^2,\vartheta \sim \mathcal N\!\big(\mu,\; v_i(\vartheta)\big),
\qquad
v_i(\vartheta)=\sigma^2\,L_i^{\beta}\,s(L_i;\xi),
\end{equation}
where \(\mu\in\mathbb R\) is the common location, \(\sigma^2>0\) is a global scale,
\(\beta\in\mathbb R\) (the length exponent) controls the length–variance law,
and \(s(L;\xi)>0\) captures within-trajectory dependence via a low-dimensional shape \(\xi\).
We denote \(\vartheta=(\beta,\xi)\).
Define the precisions and normalized weights
\begin{equation}
\label{eq:prec-weights}
\omega_i(\beta,\xi)=\frac{1}{L_i^{\beta}\,s(L_i;\xi)},
\qquad
w_i(\beta,\xi)=\frac{\omega_i(\beta,\xi)}{\sum_{j=1}^G \omega_j(\beta,\xi)}.
\end{equation}
We take \(p(\mu)\propto 1\) and Jeffreys \(p(\sigma^2)\propto 1/\sigma^2\); weak priors \(\pi_{\beta},\pi_{\xi}\) may be used or set flat for EB/type-II ML.

\paragraph{Weighted aggregator.}
The Bayes estimator of \(\mu\) under \eqref{eq:model} (for fixed \(\vartheta\)) is the precision-weighted mean
\begin{equation}
\label{eq:mu-agg}
m(\beta,\xi)=\sum_{i=1}^G w_i(\beta,\xi)\,g_i
=\frac{\sum_{i=1}^G \omega_i(\beta,\xi)\,g_i}{\sum_{i=1}^G \omega_i(\beta,\xi)}.
\end{equation}
Let
\[
e_i(\beta,\xi)=g_i-m(\beta,\xi),
\qquad
\Lambda_\omega(\beta,\xi)=\sum_{i=1}^G \omega_i(\beta,\xi),
\]
and
\begin{equation}
\label{eq:Lambda-RSS}
\mathrm{RSS}_\omega(\beta,\xi)=\sum_{i=1}^G \omega_i(\beta,\xi)\,e_i(\beta,\xi)^2.
\end{equation}

\begin{proposition}[Posterior in closed form, conditional on \((\beta,\xi)\)]
\label{prop:posterior}
Under \eqref{eq:model} with \(p(\mu)\propto 1\) and \(p(\sigma^2)\propto \frac{1}{\sigma^2}\),
\[
\mu \mid \sigma^2,\mathbf g,\vartheta \sim \mathcal N\!\big(m(\beta,\xi),\,\sigma^2/\Lambda_\omega(\beta,\xi)\big),
\quad
\sigma^2 \mid \mathbf g,\vartheta \sim \mathrm{Inv\!-\!Gamma}\!\left(\frac{G-1}{2},\,\frac{\mathrm{RSS}_\omega(\beta,\xi)}{2}\right),
\]
where \(\mathbf g=(g_1,\dots,g_G)\) and we use the inverse-gamma parameterization with density
\(p(x)\propto x^{-(a+1)}\exp(-b/x)\) for \(x>0\), shape \(a>0\) and scale \(b>0\).
\end{proposition}
\begin{proof}
We first write down the likelihood explicitly.
Given \(\vartheta=(\beta,\xi)\),
\[
v_i(\vartheta)
=\sigma^2 L_i^{\beta} s(L_i;\xi)
=\frac{\sigma^2}{\omega_i(\beta,\xi)}.
\]
Thus the conditional density of \(g_i\) is
\[
p(g_i\mid \mu,\sigma^2,\vartheta)
=\frac{1}{\sqrt{2\pi v_i(\vartheta)}}
\exp\!\left(-\frac{(g_i-\mu)^2}{2v_i(\vartheta)}\right)
=\frac{\sqrt{\omega_i}}{\sqrt{2\pi\sigma^2}}
\exp\!\left(-\frac{\omega_i}{2\sigma^2}(g_i-\mu)^2\right),
\]
where, for brevity, we write \(\omega_i=\omega_i(\beta,\xi)\).
By independence across \(i\),
\[
p(\mathbf g\mid \mu,\sigma^2,\vartheta)
=\prod_{i=1}^G p(g_i\mid \mu,\sigma^2,\vartheta)
=(2\pi\sigma^2)^{-G/2}\Bigg(\prod_{i=1}^G \omega_i^{1/2}\Bigg)
\exp\!\left(-\frac{1}{2\sigma^2}\sum_{i=1}^G \omega_i(g_i-\mu)^2\right).
\]
Multiplying by the priors \(p(\mu)\propto 1\) and \(p(\sigma^2)\propto 1/\sigma^2\) yields the unnormalized joint posterior
\begin{equation}
\label{eq:joint-posterior}
p(\mu,\sigma^2\mid\mathbf g,\vartheta)
\;\propto\;
(\sigma^2)^{-G/2}\frac{1}{\sigma^2}
\Bigg(\prod_{i=1}^G \omega_i^{1/2}\Bigg)
\exp\!\left(-\frac{1}{2\sigma^2}\sum_{i=1}^G \omega_i(g_i-\mu)^2\right).
\end{equation}
Next we complete the square in \(\mu\).
Let
\[
T=\sum_{i=1}^G \omega_i g_i,\qquad
\Lambda_\omega=\sum_{i=1}^G \omega_i,
\qquad
m=\frac{T}{\Lambda_\omega}.
\]
Then
\[
\sum_{i=1}^G \omega_i(g_i-\mu)^2
=\sum_{i=1}^G \omega_i\big[(g_i-m)+(m-\mu)\big]^2
=\sum_{i=1}^G \omega_i(g_i-m)^2
+2(m-\mu)\sum_{i=1}^G \omega_i(g_i-m)
+\Lambda_\omega(\mu-m)^2.
\]
The middle term vanishes because
\[
\sum_{i=1}^G \omega_i(g_i-m)
=\sum_{i=1}^G \omega_i g_i - m\sum_{i=1}^G \omega_i
=T-\Lambda_\omega m
=0.
\]
Hence
\[
\sum_{i=1}^G \omega_i(g_i-\mu)^2
=\Lambda_\omega(\mu-m)^2 + \sum_{i=1}^G \omega_i(g_i-m)^2
=\Lambda_\omega(\mu-m)^2 + \mathrm{RSS}_\omega,
\]
with \(\mathrm{RSS}_\omega\) as in \eqref{eq:Lambda-RSS}.
Substituting into \eqref{eq:joint-posterior} gives
\[
p(\mu,\sigma^2\mid\mathbf g,\vartheta)
\;\propto\;
(\sigma^2)^{-(G/2+1)}
\Bigg(\prod_{i=1}^G \omega_i^{1/2}\Bigg)
\exp\!\left(-\frac{\Lambda_\omega}{2\sigma^2}(\mu-m)^2\right)
\exp\!\left(-\frac{\mathrm{RSS}_\omega}{2\sigma^2}\right).
\]

\emph{Conditional posterior of \(\mu\).}
For fixed \(\sigma^2\), the dependence on \(\mu\) is exactly that of a Normal kernel with mean \(m\) and variance \(\sigma^2/\Lambda_\omega\).
Normalizing in \(\mu\) therefore yields
\[
\mu\mid\sigma^2,\mathbf g,\vartheta
\sim \mathcal N\!\big(m(\beta,\xi),\,\sigma^2/\Lambda_\omega(\beta,\xi)\big).
\]

\emph{Marginal posterior of \(\sigma^2\).}
Integrate out \(\mu\) from \(p(\mu,\sigma^2\mid\mathbf g,\vartheta)\).
The \(\mu\)-integral is
\[
\int_{-\infty}^{\infty}
\exp\!\left(-\frac{\Lambda_\omega}{2\sigma^2}(\mu-m)^2\right)\,d\mu
=\sqrt{\frac{2\pi\sigma^2}{\Lambda_\omega}}.
\]
Thus
\[
p(\sigma^2\mid\mathbf g,\vartheta)
\;\propto\;
(\sigma^2)^{-(G/2+1)}
\sqrt{\frac{2\pi\sigma^2}{\Lambda_\omega}}
\Bigg(\prod_{i=1}^G \omega_i^{1/2}\Bigg)
\exp\!\left(-\frac{\mathrm{RSS}_\omega}{2\sigma^2}\right).
\]
The multiplicative constant
\(\sqrt{2\pi/\Lambda_\omega}\prod_i \omega_i^{1/2}\) does not depend on \(\sigma^2\), so we drop it.
We are left with
\[
p(\sigma^2\mid\mathbf g,\vartheta)
\;\propto\;
(\sigma^2)^{-(G/2+1-1/2)}\exp\!\left(-\frac{\mathrm{RSS}_\omega}{2\sigma^2}\right)
=(\sigma^2)^{-(G+1)/2}\exp\!\left(-\frac{\mathrm{RSS}_\omega}{2\sigma^2}\right).
\]
Comparing with the inverse-gamma kernel
\(x^{-(a+1)}\exp(-b/x)\) on \(x>0\),
we read off \(a=(G-1)/2\) and \(b=\mathrm{RSS}_\omega/2\).
This proves the claimed inverse-gamma posterior for \(\sigma^2\).
\end{proof}

\begin{proposition}[GLS optimality of the aggregator]
\label{prop:gls}
Fix \((\beta,\xi)\).
Among all linear unbiased estimators \(\tilde\mu=\sum_{i=1}^G a_i g_i\) with \(\sum_{i=1}^G a_i=1\),
\(m(\beta,\xi)\) in \eqref{eq:mu-agg} attains the minimum variance
\(\mathrm{Var}(\tilde\mu)=\sigma^2/\Lambda_\omega(\beta,\xi)\).
\end{proposition}
\begin{proof}
For fixed \(\vartheta\), the model \eqref{eq:model} implies
\[
\mathbb E[g_i\mid\mu,\vartheta]=\mu,
\qquad
\mathrm{Var}(g_i\mid\mu,\vartheta)=v_i(\vartheta)
=\frac{\sigma^2}{\omega_i},
\]
and the \(g_i\) are conditionally independent.
Consider a generic linear estimator
\(\tilde\mu=\sum_i a_i g_i\) with coefficients \(a_i\in\mathbb R\).

\emph{Unbiasedness constraint.}
\[
\mathbb E[\tilde\mu\mid\mu,\vartheta]
=\sum_{i=1}^G a_i \mathbb E[g_i\mid\mu,\vartheta]
=\sum_{i=1}^G a_i \mu
=\mu \sum_{i=1}^G a_i.
\]
Thus \(\tilde\mu\) is unbiased for \(\mu\) if and only if
\(\sum_{i=1}^G a_i=1\).

\emph{Variance under the constraint.}
Using independence and \(\mathrm{Var}(g_i)=\sigma^2/\omega_i\),
\[
\mathrm{Var}(\tilde\mu\mid\mu,\vartheta)
=\sum_{i=1}^G a_i^2\,\mathrm{Var}(g_i\mid\mu,\vartheta)
=\sigma^2 \sum_{i=1}^G \frac{a_i^2}{\omega_i}.
\]
We must minimize \(\sum_i a_i^2/\omega_i\) subject to \(\sum_i a_i=1\).
Introduce a Lagrange multiplier \(\lambda\) and consider
\[
\mathcal L(a_1,\dots,a_G,\lambda)
=\sum_{i=1}^G \frac{a_i^2}{\omega_i}
-\lambda\Bigg(\sum_{i=1}^G a_i -1\Bigg).
\]
Setting partial derivatives to zero,
\[
\frac{\partial \mathcal L}{\partial a_i}
=\frac{2a_i}{\omega_i}-\lambda
=0
\quad\Rightarrow\quad
a_i=\frac{\lambda}{2}\,\omega_i,\qquad i=1,\dots,G.
\]
Imposing \(\sum_i a_i=1\) gives
\[
\sum_{i=1}^G a_i
=\frac{\lambda}{2}\sum_{i=1}^G \omega_i
=\frac{\lambda}{2}\Lambda_\omega
=1
\quad\Rightarrow\quad
\lambda=\frac{2}{\Lambda_\omega}.
\]
Thus
\[
a_i
=\frac{\omega_i}{\Lambda_\omega}
=w_i(\beta,\xi),
\]
and the unique minimizer is exactly the weight vector in \eqref{eq:mu-agg}.
The corresponding variance is
\[
\mathrm{Var}(m(\beta,\xi)\mid\mu,\vartheta)
=\sigma^2 \sum_{i=1}^G \frac{w_i^2}{\omega_i}
=\sigma^2 \sum_{i=1}^G \frac{\omega_i^2}{\Lambda_\omega^2\,\omega_i}
=\sigma^2 \frac{\sum_i \omega_i}{\Lambda_\omega^2}
=\frac{\sigma^2}{\Lambda_\omega}.
\]
Hence \(m(\beta,\xi)\) is the Gauss–Markov / GLS estimator with minimum variance among all linear unbiased estimators.
\end{proof}

\begin{theorem}[EB marginal objective for \((\beta,\xi)\)]
\label{thm:EB-obj}
Under \eqref{eq:model} with the above priors, the marginal type-II log-posterior for \(\vartheta=(\beta,\xi)\) is (up to an additive constant)
\begin{equation}
\label{eq:EB-obj}
\ell(\beta,\xi)
=\log \pi_\beta(\beta)+\log \pi_\xi(\xi)
+\tfrac{1}{2}\sum_{i=1}^G \log \omega_i(\beta,\xi)
-\tfrac{1}{2}\log \Lambda_\omega(\beta,\xi)
-\tfrac{G-1}{2}\log \mathrm{RSS}_\omega(\beta,\xi).
\end{equation}
\end{theorem}
\begin{proof}
We integrate out \(\mu\) and \(\sigma^2\) from the joint model
\[
p(\mathbf g,\mu,\sigma^2\mid\vartheta)
=p(\mathbf g\mid\mu,\sigma^2,\vartheta)\,p(\mu)\,p(\sigma^2)
\]
and then add the priors \(\pi_\beta,\pi_\xi\) on \((\beta,\xi)\).

From the likelihood calculation in the proof of Proposition~\ref{prop:posterior}, we have
\[
p(\mathbf g\mid \mu,\sigma^2,\vartheta)
=(2\pi\sigma^2)^{-G/2}\left(\prod_{i=1}^G \omega_i^{1/2}\right)
\exp\!\left(-\frac{1}{2\sigma^2}\sum_{i=1}^G \omega_i(g_i-\mu)^2\right),
\]
and
\(\sum_{i=1}^G \omega_i(g_i-\mu)^2
=\Lambda_\omega(\mu-m)^2+\mathrm{RSS}_\omega\).
Using \(p(\mu)\propto 1\) and \(p(\sigma^2)\propto 1/\sigma^2\),
\[
p(\mathbf g,\mu,\sigma^2\mid\vartheta)
\propto
(\sigma^2)^{-(G/2+1)}
\left(\prod_{i=1}^G \omega_i^{1/2}\right)
\exp\!\left(-\frac{\Lambda_\omega}{2\sigma^2}(\mu-m)^2\right)
\exp\!\left(-\frac{\mathrm{RSS}_\omega}{2\sigma^2}\right).
\]
Integrate with respect to \(\mu\).
As in Proposition~\ref{prop:posterior},
\[
\int_{-\infty}^\infty
\exp\!\left(-\frac{\Lambda_\omega}{2\sigma^2}(\mu-m)^2\right)\,d\mu
=\sqrt{\frac{2\pi\sigma^2}{\Lambda_\omega}},
\]
hence
\[
p(\mathbf g,\sigma^2\mid\vartheta)
\propto
(\sigma^2)^{-(G/2+1)}
\sqrt{\frac{2\pi\sigma^2}{\Lambda_\omega}}
\left(\prod_{i=1}^G \omega_i^{1/2}\right)
\exp\!\left(-\frac{\mathrm{RSS}_\omega}{2\sigma^2}\right)
\]
\[
\propto
(\sigma^2)^{-(G+1)/2}
\left(\prod_{i=1}^G \omega_i^{1/2}\right)
\Lambda_\omega^{-1/2}
\exp\!\left(-\frac{\mathrm{RSS}_\omega}{2\sigma^2}\right),
\]
where we have absorbed constants depending on \(2\pi\) into a factor that does not depend on \(\vartheta\).

Now integrate out \(\sigma^2\) using the standard identity for inverse-gamma integrals:
\[
\int_0^\infty
(\sigma^2)^{-(a+1)} \exp\!\left(-\frac{b}{\sigma^2}\right)\,d\sigma^2
=2\,b^{-a}\,\Gamma(a),
\qquad a>0,\,b>0.
\]
Here we take
\(a=(G-1)/2\), \(b=\mathrm{RSS}_\omega/2\).
Thus
\[
\int_0^\infty
(\sigma^2)^{-(G+1)/2}
\exp\!\left(-\frac{\mathrm{RSS}_\omega}{2\sigma^2}\right)\,d\sigma^2
\propto \mathrm{RSS}_\omega^{-(G-1)/2}.
\]
It follows that the marginal likelihood of the data given \(\vartheta\) is
\[
p(\mathbf g\mid\vartheta)
\propto
\left(\prod_{i=1}^G \omega_i^{1/2}\right)
\Lambda_\omega^{-1/2}
\mathrm{RSS}_\omega^{-(G-1)/2},
\]
up to factors independent of \(\vartheta\).
Taking logarithms and adding \(\log \pi_\beta(\beta)+\log \pi_\xi(\xi)\) yields
% \[
% \ell(\beta,\xi)
% =\log p(\mathbf g\mid\vartheta)
% +\log \pi_\beta(\beta)+\log \pi_\xi(\xi)
% =\log \pi_\beta(\beta)+\log \pi_\xi(\xi)
% +\frac{1}{2}\sum_{i=1}^G \log \omega_i
% -\frac{1}{2}\log \Lambda_\omega
% -\frac{G-1}{2}\log \mathrm{RSS}_\omega
% + C,
% \]

\begin{align}
\ell(\beta,\xi)
&= \log p(\mathbf g \mid \vartheta)
  + \log \pi_\beta(\beta)
  + \log \pi_\xi(\xi) \\
&= \log \pi_\beta(\beta)
  + \log \pi_\xi(\xi)
  + \frac{1}{2}\sum_{i=1}^G \log \omega_i
  - \frac{1}{2}\log \Lambda_\omega \\
&\quad
  - \frac{G-1}{2}\log \mathrm{RSS}_\omega
  + C .
\end{align}
where \(C\) does not depend on \((\beta,\xi)\).
Dropping the additive constant proves \eqref{eq:EB-obj}.
\end{proof}

\subsubsection{Prompt-group (GRPO) extension: group-specific means}\label{subsec:eb-grouped}
In GRPO-style training, rollouts are naturally organized into \emph{prompt groups}: for each prompt $p\in\{1,\dots,P\}$ we sample $N_p\ge 2$ trajectories and construct within-group normalized advantages/baselines. To align the EB variance-law estimation with this grouping while preserving the closed-form structure above, we extend \eqref{eq:model} to allow a \emph{group-specific} location parameter $\mu_p$ that is integrated out with the same flat prior.

Let $\mathcal I_p$ be the index set of trajectories in group $p$, with $|\mathcal I_p|=N_p$ and total $G:=\sum_{p=1}^P N_p$. We observe $(g_{p,i},L_{p,i})$ for $i\in\mathcal I_p$, and adopt the grouped working model
\begin{equation}
\label{eq:model-grouped}
g_{p,i}\mid \mu_p,\sigma^2,\vartheta \sim \mathcal N\!\Big(\mu_p,\; v_{p,i}(\vartheta)\Big),\qquad
v_{p,i}(\vartheta)=\sigma^2 L_{p,i}^{\beta}s(L_{p,i};\xi),
\end{equation}
with independent improper priors $p(\mu_p)\propto 1$ and $p(\sigma^2)\propto 1/\sigma^2$.
Define group precisions and weights
\begin{equation}
\label{eq:prec-weights-grouped}
\omega_{p,i}(\beta,\xi)=\frac{1}{L_{p,i}^{\beta}s(L_{p,i};\xi)},\qquad
\Lambda_{\omega,p}=\sum_{j\in\mathcal I_p}\omega_{p,j},\qquad
w_{p,i}=\frac{\omega_{p,i}}{\Lambda_{\omega,p}},
\end{equation}
and the within-group precision-weighted aggregator
\begin{equation}
\label{eq:mu-agg-grouped}
m_p(\beta,\xi)=\sum_{i\in\mathcal I_p} w_{p,i}(\beta,\xi)\,g_{p,i}.
\end{equation}
Let $e_{p,i}(\beta,\xi)=g_{p,i}-m_p(\beta,\xi)$ and define the grouped RSS
\begin{equation}
\label{eq:RSS-grouped}
\mathrm{RSS}_{\omega}^{\mathrm{grp}}(\beta,\xi)=\sum_{p=1}^P\sum_{i\in\mathcal I_p}\omega_{p,i}(\beta,\xi)\,e_{p,i}(\beta,\xi)^2.
\end{equation}

\begin{theorem}[Grouped EB marginal objective for $\,\vartheta=(\beta,\xi)$]
\label{thm:EB-obj-grouped}
Under \eqref{eq:model-grouped} with the above priors, the marginal type-II log-posterior for $\vartheta=(\beta,\xi)$ is (up to an additive constant)
\begin{equation}
\label{eq:EB-obj-grouped}
\ell_{\mathrm{grp}}(\beta,\xi)
=\log \pi_\beta(\beta)+\log \pi_\xi(\xi)
+\tfrac{1}{2}\sum_{p=1}^P\sum_{i\in\mathcal I_p}\log \omega_{p,i}(\beta,\xi)
-\tfrac{1}{2}\sum_{p=1}^P\log \Lambda_{\omega,p}(\beta,\xi)
-\tfrac{G-P}{2}\log \mathrm{RSS}_{\omega}^{\mathrm{grp}}(\beta,\xi).
\end{equation}
When $P=1$ (a single group), \eqref{eq:EB-obj-grouped} reduces to \eqref{eq:EB-obj}.
\end{theorem}
\begin{proof}
Conditioned on $\vartheta$, the likelihood factorizes over groups:
$p(\{g_{p,i}\}\mid\{\mu_p\},\sigma^2,\vartheta)=\prod_{p=1}^P\prod_{i\in\mathcal I_p} \mathcal N\!\big(g_{p,i};\mu_p,\sigma^2/\omega_{p,i}\big)$.
For each group $p$, completing the square yields
\(\sum_{i\in\mathcal I_p}\omega_{p,i}(g_{p,i}-\mu_p)^2=\Lambda_{\omega,p}(\mu_p-m_p)^2+\sum_{i\in\mathcal I_p}\omega_{p,i}(g_{p,i}-m_p)^2\).
Integrating out $\mu_p$ under $p(\mu_p)\propto 1$ gives a factor $\sqrt{2\pi\sigma^2/\Lambda_{\omega,p}}$ per group, hence the grouped marginal (conditional on $\sigma^2$) is proportional to
\((\sigma^2)^{-G/2}\prod_{p,i}\omega_{p,i}^{1/2}\cdot \prod_{p}\Lambda_{\omega,p}^{-1/2}\cdot\exp\{-\mathrm{RSS}_{\omega}^{\mathrm{grp}}/(2\sigma^2)\}\).
Multiplying by $p(\sigma^2)\propto 1/\sigma^2$ and integrating out $\sigma^2$ produces the exponent $(G-P)/2$ in \eqref{eq:EB-obj-grouped}, since integrating out $P$ group means reduces the degrees of freedom by $P$.
Taking logarithms and adding $\log\pi_\beta+\log\pi_\xi$ yields \eqref{eq:EB-obj-grouped}.
\end{proof}

\begin{lemma}[Weighted RSS derivative identity]
\label{lem:rss-deriv}
$\forall$ component \(\phi\in\{\beta\}\cup\{\xi\text{ components}\}\),
\begin{equation}
\label{eq:rss-deriv}
\frac{\partial}{\partial \phi}\,\mathrm{RSS}_\omega(\beta,\xi)
=\sum_{i=1}^G \frac{\partial \omega_i}{\partial \phi}\,e_i^2
=\sum_{i=1}^G \omega_i\,e_i^2\,\frac{\partial \log \omega_i}{\partial \phi}.
\end{equation}
\end{lemma}
\begin{proof}
Recall
\(\mathrm{RSS}_\omega(\beta,\xi)=\sum_{i=1}^G \omega_i e_i^2\),
with \(e_i=g_i-m\) and \(m=m(\beta,\xi)\).
We differentiate \(\mathrm{RSS}_\omega\) with respect to \(\phi\).
Using the product and chain rules,
\[
\frac{\partial}{\partial\phi}\,\mathrm{RSS}_\omega
=\sum_{i=1}^G \left[
\frac{\partial \omega_i}{\partial\phi}\,e_i^2
+\omega_i \frac{\partial}{\partial\phi}(e_i^2)
\right].
\]
Since \(e_i=g_i-m\) and \(g_i\) does not depend on \(\phi\),
\(\frac{\partial e_i}{\partial\phi}=-\frac{\partial m}{\partial\phi}\), so
\[
\frac{\partial}{\partial\phi}(e_i^2)
=2e_i\,\frac{\partial e_i}{\partial\phi}
=2e_i\left(-\frac{\partial m}{\partial\phi}\right)
=-2e_i\,\frac{\partial m}{\partial\phi}.
\]
Therefore,
\[
\frac{\partial}{\partial\phi}\,\mathrm{RSS}_\omega
=\sum_{i=1}^G \frac{\partial \omega_i}{\partial\phi}\,e_i^2
-2\frac{\partial m}{\partial\phi}\sum_{i=1}^G \omega_i e_i.
\]
To simplify the second term, observe that
\[
\sum_{i=1}^G \omega_i e_i
=\sum_{i=1}^G \omega_i(g_i-m)
=\sum_{i=1}^G \omega_i g_i - m\sum_{i=1}^G \omega_i
=T - m\Lambda_\omega
=0,
\]
because \(m=T/\Lambda_\omega\).
Thus the term involving \(\partial m/\partial\phi\) cancels, and we obtain
\[
\frac{\partial}{\partial\phi}\,\mathrm{RSS}_\omega
=\sum_{i=1}^G \frac{\partial \omega_i}{\partial\phi}\,e_i^2.
\]
Finally, use \(\frac{\partial \omega_i}{\partial\phi}=\omega_i\,\frac{\partial \log \omega_i}{\partial\phi}\) to get the second equality:
\[
\sum_{i=1}^G \frac{\partial \omega_i}{\partial\phi}\,e_i^2
=\sum_{i=1}^G \omega_i e_i^2\,\frac{\partial \log \omega_i}{\partial\phi}.
\]
\end{proof}

\begin{proposition}[Gradient of the EB objective]
\label{prop:grad}
With \(\Lambda_\omega,\mathrm{RSS}_\omega\) as above,
\begin{equation}
\label{eq:EB-grad-general}
\frac{\partial \ell}{\partial \phi}
=\frac{1}{2}\!\left[
\sum_{i=1}^G \frac{\partial \log \omega_i}{\partial \phi}
-\frac{1}{\Lambda_\omega}\sum_{i=1}^G \frac{\partial \omega_i}{\partial \phi}
\right]
-\frac{G-1}{2}\frac{1}{\mathrm{RSS}_\omega}
\sum_{i=1}^G \omega_i e_i^2\,\frac{\partial \log \omega_i}{\partial \phi}
+\frac{\partial \log \pi(\phi)}{\partial \phi},
\end{equation}
where \(\pi(\phi)\) denotes the marginal prior factor in \(\pi_\beta\pi_\xi\) involving \(\phi\).
\end{proposition}
\begin{proof}
Starting from \eqref{eq:EB-obj}, we differentiate term by term with respect to \(\phi\).

\emph{Prior term.}
By definition,
\[
\frac{\partial}{\partial\phi}\big(\log \pi_\beta(\beta)+\log \pi_\xi(\xi)\big)
=\frac{\partial \log \pi(\phi)}{\partial \phi},
\]
where \(\pi(\phi)\) denotes the appropriate prior factor.

\emph{Sum of log-precisions.}
For the third term in \eqref{eq:EB-obj},
\[
\frac{\partial}{\partial\phi}\left(\frac{1}{2}\sum_{i=1}^G \log\omega_i\right)
=\frac{1}{2}\sum_{i=1}^G \frac{1}{\omega_i}\frac{\partial \omega_i}{\partial\phi}
=\frac{1}{2}\sum_{i=1}^G \frac{\partial \log \omega_i}{\partial\phi}.
\]

\emph{Log of total precision.}
For the fourth term,
\[
\frac{\partial}{\partial\phi}\left(-\frac{1}{2}\log\Lambda_\omega\right)
=-\frac{1}{2}\frac{1}{\Lambda_\omega}\frac{\partial \Lambda_\omega}{\partial\phi}
=-\frac{1}{2}\frac{1}{\Lambda_\omega}\sum_{i=1}^G \frac{\partial \omega_i}{\partial\phi}.
\]

\emph{Residual sum of squares term.}
For the final term,
\[
\frac{\partial}{\partial\phi}\left(-\frac{G-1}{2}\log \mathrm{RSS}_\omega\right)
=-\frac{G-1}{2}\frac{1}{\mathrm{RSS}_\omega}
\frac{\partial \mathrm{RSS}_\omega}{\partial\phi}.
\]
By Lemma~\ref{lem:rss-deriv},
\(\frac{\partial \mathrm{RSS}_\omega}{\partial\phi}
=\sum_i \omega_i e_i^2\,\frac{\partial \log \omega_i}{\partial\phi}\),
hence
\[
-\frac{G-1}{2}\frac{1}{\mathrm{RSS}_\omega}
\frac{\partial \mathrm{RSS}_\omega}{\partial\phi}
=-\frac{G-1}{2}\frac{1}{\mathrm{RSS}_\omega}
\sum_{i=1}^G \omega_i e_i^2\,\frac{\partial \log \omega_i}{\partial\phi}.
\]

Collecting the contributions from all terms gives \eqref{eq:EB-grad-general}.
\end{proof}

\begin{corollary}[Gradient in the length direction]
\label{cor:grad-beta}
Since \(\log\omega_i=-\beta\log L_i-\log s(L_i;\xi)\),
\(\partial \log\omega_i/\partial\beta=-\log L_i\) and \(\partial \omega_i/\partial\beta=-\omega_i \log L_i\).
Hence
\begin{equation}
\label{eq:EB-grad-beta}
\frac{\partial \ell}{\partial \beta}
=\frac{1}{2}\!\left[-\sum_{i=1}^G \log L_i
+\frac{\sum_{i=1}^G \omega_i \log L_i}{\Lambda_\omega}\right]
+\frac{G-1}{2}\frac{1}{\mathrm{RSS}_\omega}\sum_{i=1}^G \omega_i e_i^2 \log L_i
+\frac{\partial \log \pi_\beta(\beta)}{\partial \beta}.
\end{equation}
\end{corollary}
\begin{proof}
Apply Proposition~\ref{prop:grad} with \(\phi=\beta\).
We have
\[
\frac{\partial \log\omega_i}{\partial\beta}
=-\log L_i,
\qquad
\frac{\partial \omega_i}{\partial\beta}
=\omega_i\,\frac{\partial \log\omega_i}{\partial\beta}
=-\omega_i \log L_i.
\]
Substitute these into \eqref{eq:EB-grad-general}.
The first bracketed term becomes
\[
\frac{1}{2}\left[
\sum_{i=1}^G \frac{\partial \log\omega_i}{\partial\beta}
-\frac{1}{\Lambda_\omega}\sum_{i=1}^G \frac{\partial \omega_i}{\partial\beta}
\right]
=\frac{1}{2}\left[
\sum_{i=1}^G (-\log L_i)
-\frac{1}{\Lambda_\omega}\sum_{i=1}^G (-\omega_i \log L_i)
\right]
\]
\[
=\frac{1}{2}\left[
-\sum_{i=1}^G \log L_i
+\frac{\sum_{i=1}^G \omega_i \log L_i}{\Lambda_\omega}
\right].
\]
The second term in \eqref{eq:EB-grad-general} becomes
\[
-\frac{G-1}{2}\frac{1}{\mathrm{RSS}_\omega}
\sum_{i=1}^G \omega_i e_i^2\,\frac{\partial \log\omega_i}{\partial\beta}
=-\frac{G-1}{2}\frac{1}{\mathrm{RSS}_\omega}
\sum_{i=1}^G \omega_i e_i^2 (-\log L_i)
\]
\[
=\frac{G-1}{2}\frac{1}{\mathrm{RSS}_\omega}
\sum_{i=1}^G \omega_i e_i^2 \log L_i.
\]
Adding the prior derivative \(\partial \log \pi_\beta(\beta)/\partial\beta\) yields \eqref{eq:EB-grad-beta}.
\end{proof}

\begin{corollary}[Grouped gradient in the length direction (prompt groups)]
\label{cor:grad-beta-grouped}
Under the grouped EB objective \eqref{eq:EB-obj-grouped}, the gradient with respect to the length exponent $\beta$ is
\begin{equation}
\label{eq:EB-grad-beta-grouped}
\frac{\partial \ell_{\mathrm{grp}}}{\partial \beta}
=\frac{1}{2}\!\left[-\sum_{p=1}^P\sum_{i\in\mathcal I_p} \log L_{p,i}
+\sum_{p=1}^P\frac{\sum_{i\in\mathcal I_p} \omega_{p,i} \log L_{p,i}}{\Lambda_{\omega,p}}\right]
+\frac{G-P}{2}\frac{1}{\mathrm{RSS}_{\omega}^{\mathrm{grp}}}\sum_{p=1}^P\sum_{i\in\mathcal I_p} \omega_{p,i} e_{p,i}^2 \log L_{p,i}
+\frac{\partial \log \pi_\beta(\beta)}{\partial \beta}.
\end{equation}
\end{corollary}
\begin{proof}
Differentiate \eqref{eq:EB-obj-grouped} with respect to $\beta$, using $\partial_\beta\log \omega_{p,i}=-\log L_{p,i}$ and $\partial_\beta \omega_{p,i}=-\omega_{p,i}\log L_{p,i}$. The second term in \eqref{eq:EB-obj-grouped} contributes a sum of groupwise normalizers $\Lambda_{\omega,p}$, and integrating out $P$ group means yields the replacement $G-1\mapsto G-P$ in the RSS exponent. The resulting expression is \eqref{eq:EB-grad-beta-grouped}.
\end{proof}

\begin{definition}[Stretched-geometric \(k\)-band dependence]
\label{def:kband}
Given \(k\in\mathbb N\), \(\rho\in(-1,1)\), \(\eta\ge 0\), define with \(m=\min\{k,L-1\}\)
\begin{equation}
\label{eq:s-family}
s(L;\rho,k,\eta)=1+\frac{2}{L}\sum_{h=1}^m (L-h)\,\rho^{\,h^{\eta}}.
\end{equation}
\end{definition}

\begin{lemma}[Derivatives of the \(k\)-band shape]
\label{lem:kband-deriv}
Let \(s_i=s(L_i;\rho,k,\eta)\), \(m_i=\min\{k,L_i-1\}\).
Then
\[
\frac{\partial \log \omega_i}{\partial \rho}
=-\frac{1}{s_i}\frac{\partial s_i}{\partial \rho},
\quad
\frac{\partial s_i}{\partial \rho}
=\frac{2}{L_i}\sum_{h=1}^{m_i} (L_i-h)\, h^{\eta}\, \rho^{\,h^{\eta}-1},
\]
\[
\frac{\partial \log \omega_i}{\partial \eta}
=-\frac{1}{s_i}\frac{\partial s_i}{\partial \eta},
\quad
\frac{\partial s_i}{\partial \eta}
=\frac{2}{L_i}\sum_{h=1}^{m_i} (L_i-h)\, \rho^{\,h^{\eta}} (\log \rho)\, h^{\eta}\log h.
\]
\end{lemma}
\begin{proof}
By Definition~\ref{def:kband},
\[
s_i
=1+\frac{2}{L_i}\sum_{h=1}^{m_i} (L_i-h)\,\rho^{\,h^{\eta}}.
\]
For \(\partial s_i/\partial \rho\), note that \(\partial(1)/\partial\rho=0\) and
we can differentiate the sum termwise:
\[
\frac{\partial s_i}{\partial \rho}
=\frac{2}{L_i}\sum_{h=1}^{m_i} (L_i-h)\,\frac{\partial}{\partial\rho}\big(\rho^{\,h^{\eta}}\big).
\]
Using \(\rho^{\,h^{\eta}}=\exp(h^{\eta}\log\rho)\), the chain rule gives
\[
\frac{\partial}{\partial\rho}\big(\rho^{\,h^{\eta}}\big)
=h^{\eta} \rho^{\,h^{\eta}-1},
\]
so
\[
\frac{\partial s_i}{\partial \rho}
=\frac{2}{L_i}\sum_{h=1}^{m_i} (L_i-h)\,h^{\eta}\,\rho^{\,h^{\eta}-1}.
\]

For \(\partial s_i/\partial \eta\), again differentiate termwise:
\[
\frac{\partial s_i}{\partial \eta}
=\frac{2}{L_i}\sum_{h=1}^{m_i} (L_i-h)\,\frac{\partial}{\partial\eta}\big(\rho^{\,h^{\eta}}\big).
\]
Write \(\rho^{\,h^{\eta}}=\exp(h^{\eta}\log\rho)\).
Then
\[
\frac{\partial}{\partial\eta}\big(\rho^{\,h^{\eta}}\big)
=\rho^{\,h^{\eta}}\frac{\partial}{\partial\eta}\big(h^{\eta}\log\rho\big)
=\rho^{\,h^{\eta}}(\log\rho)h^{\eta}\log h.
\]
Thus
\[
\frac{\partial s_i}{\partial \eta}
=\frac{2}{L_i}\sum_{h=1}^{m_i} (L_i-h)\,\rho^{\,h^{\eta}}(\log\rho)h^{\eta}\log h.
\]

Finally, since \(\omega_i=L_i^{-\beta}s(L_i;\xi)^{-1}\),
\[
\log\omega_i
=-\beta\log L_i-\log s_i,
\]
so
\[
\frac{\partial \log \omega_i}{\partial\rho}
=-\frac{1}{s_i}\frac{\partial s_i}{\partial\rho},
\qquad
\frac{\partial \log \omega_i}{\partial\eta}
=-\frac{1}{s_i}\frac{\partial s_i}{\partial\eta}.
\]
\end{proof}

\begin{proposition}[Unbiasedness of \(m(\beta,\xi)\)]
\label{prop:unbiased}
For fixed \((\beta,\xi)\), the estimator \(m(\beta,\xi)\) in \eqref{eq:mu-agg} is unbiased for \(\mu\).
\end{proposition}
\begin{proof}
By \eqref{eq:model}, for each \(i\),
\(\mathbb E[g_i\mid\mu,\vartheta]=\mu\).
Using \eqref{eq:mu-agg},
\[
\mathbb E[m(\beta,\xi)\mid\mu,\vartheta]
=\mathbb E\Bigg[\sum_{i=1}^G w_i g_i\;\bigg|\;\mu,\vartheta\Bigg]
=\sum_{i=1}^G w_i\,\mathbb E[g_i\mid\mu,\vartheta]
=\mu\sum_{i=1}^G w_i.
\]
From \eqref{eq:prec-weights}, \(\sum_i w_i=1\).
Therefore \(\mathbb E[m(\beta,\xi)\mid\mu,\vartheta]=\mu\), i.e., \(m(\beta,\xi)\) is unbiased.
\end{proof}

\begin{proposition}[Link to the \(\Delta L\) rule]
\label{prop:deltaL}
If \(s(L;\xi)\equiv 1\) and \(\beta=1\), then \(\omega_i=1/L_i\) and \(w_i \propto 1/L_i\), recovering the \(\Delta L\) baseline weighting.
\end{proposition}
\begin{proof}
If \(s(L;\xi)\equiv 1\) and \(\beta=1\), then for each \(i\),
\[
\omega_i
=\frac{1}{L_i^{\beta}s(L_i;\xi)}
=\frac{1}{L_i\cdot 1}
=\frac{1}{L_i}.
\]
Hence
\[
w_i
=\frac{\omega_i}{\sum_{j=1}^G \omega_j}
=\frac{L_i^{-1}}{\sum_{j=1}^G L_j^{-1}}
\propto \frac{1}{L_i},
\]
which is the standard \(\Delta L\) choice of weights.
\end{proof}

\begin{assumption}[Regularity and identifiability]
\label{assump:reg}
(i) \(\{(g_i,L_i)\}_{i\ge 1}\) are i.i.d.\ or strictly stationary and ergodic with \(\sup_i \mathbb E[g_i^4]<\infty\).

(ii) The parameter space \(\Theta=\{\beta\}\times\Xi\) is compact (for example, \(\beta\in[\beta_{\min},\beta_{\max}]\), \(|\rho|\le \rho_{\max}<1\), \(\eta\in[0,\eta_{\max}]\), and fixed \(k\)).

(iii) The population objective \(\ell_\infty(\vartheta):=\lim_{G\to\infty}G^{-1}\,\mathbb E[\ell_G(\vartheta)]\) exists and has a unique maximizer \(\vartheta^\star=(\beta^\star,\xi^\star)\).

(iv) \(s(L;\xi)\) is smooth in \(\xi\) and bounded away from \(0\) and \(+\infty\) on \(\mathrm{supp}(L)\times\Xi\).
\end{assumption}

Here \(\ell_G(\vartheta)\) denotes the EB objective \eqref{eq:EB-obj} built from the first \(G\) trajectories; we may equivalently work with the rescaled version \(G^{-1}\ell_G(\vartheta)\), which only differs by a factor independent of \(\vartheta\).

\begin{theorem}[Consistency \& asymptotic normality of the EB estimator]
\label{thm:consistency}
Let \(\widehat\vartheta_G=\arg\max_{\vartheta\in\Theta}\ell_G(\vartheta)\) be any (measurable) maximizer.
Under Assumption~\ref{assump:reg}, \(\widehat\vartheta_G \xrightarrow{p} \vartheta^\star\) as \(G\to\infty\).

If, in addition, \(\ell_\infty\) is twice continuously differentiable at \(\vartheta^\star\) with negative-definite Hessian \(H=-\nabla^2\ell_\infty(\vartheta^\star)\), and
the gradient process satisfies a suitable central limit theorem (for example,
\(\sqrt{G}\,\nabla \ell_G(\vartheta^\star)\xrightarrow{d}\mathcal N(0,J)\) for some positive semi-definite matrix \(J\)),
then
\[
\sqrt{G}(\widehat\vartheta_G-\vartheta^\star)\ \xrightarrow{d}\ \mathcal N\!\big(0,\; H^{-1}\,J\,H^{-1}\big),
\]
where \(J\) is the asymptotic covariance of \(\sqrt{G}\,\nabla \ell_G(\vartheta^\star)\).
\end{theorem}
\begin{proof}
\emph{Consistency.}
We can write \(\ell_G(\vartheta)\) as a continuous function of empirical averages of bounded (or suitably integrable) functions of \((g_i,L_i)\).
Indeed, the terms \(\sum_i \log\omega_i\), \(\Lambda_\omega\), and \(\mathrm{RSS}_\omega\) are all sums of such functions, and the mapping from these empirical sums to \(\ell_G(\vartheta)\) is continuous in both the empirical sums and in \(\vartheta\).
By the (vector) law of large numbers and the continuous mapping theorem, for each fixed \(\vartheta\), the rescaled objective \(G^{-1}\ell_G(\vartheta)\) converges almost surely to a deterministic limit \(\ell_\infty(\vartheta)\).
Under the compactness of \(\Theta\) and smoothness/boundedness conditions in Assumption~\ref{assump:reg}, standard results in M-estimation (e.g., an argmax continuous mapping theorem) imply that this convergence is uniform on \(\Theta\).
By assumption, \(\ell_\infty\) has a unique maximizer \(\vartheta^\star\).
The uniform convergence combined with uniqueness yields
\(\widehat\vartheta_G\xrightarrow{p}\vartheta^\star\) (see, e.g., \citet{vandervaart1998}, Theorem~5.7).

\emph{Asymptotic normality.}
Consider the first-order optimality condition
\[
\nabla \ell_G(\widehat\vartheta_G)=0.
\]
Assume \(\ell_G\) is twice continuously differentiable in a neighborhood of \(\vartheta^\star\) with probability tending to one, and expand \(\nabla \ell_G\) in a Taylor series around \(\vartheta^\star\):
\[
\nabla \ell_G(\widehat\vartheta_G)
=\nabla \ell_G(\vartheta^\star)
+\nabla^2\ell_G(\tilde\vartheta_G)(\widehat\vartheta_G-\vartheta^\star),
\]
for some intermediate point \(\tilde\vartheta_G\) on the line segment between \(\widehat\vartheta_G\) and \(\vartheta^\star\).
Setting \(\nabla \ell_G(\widehat\vartheta_G)=0\) and rearranging gives
\[
\sqrt{G}(\widehat\vartheta_G-\vartheta^\star)
=-\Big(\nabla^2\ell_G(\tilde\vartheta_G)\Big)^{-1}\sqrt{G}\,\nabla \ell_G(\vartheta^\star).
\]
By the consistency already established, \(\tilde\vartheta_G\xrightarrow{p}\vartheta^\star\).
Under the stated differentiability and regularity assumptions,
\(\nabla^2\ell_G(\tilde\vartheta_G)\) converges in probability to \(\nabla^2\ell_\infty(\vartheta^\star)=-H\),
so \(\big(\nabla^2\ell_G(\tilde\vartheta_G)\big)^{-1}\xrightarrow{p}H^{-1}\).

By assumption,
\(\sqrt{G}\,\nabla \ell_G(\vartheta^\star)\xrightarrow{d}\mathcal N(0,J)\).
Applying Slutsky’s theorem to the product of the converging matrices and vectors yields
\[
\sqrt{G}(\widehat\vartheta_G-\vartheta^\star)
\xrightarrow{d}
-H^{-1}\,\mathcal N(0,J)
=\mathcal N(0,H^{-1}JH^{-1}),
\]
which proves the stated asymptotic normality.
\end{proof}

\begin{remark}[Smooth but not globally concave]
While \(\ell(\beta,\xi)\) is generally not globally concave, the gradient \(\nabla \ell\) and local Hessian \(\nabla^2\ell\) are explicit via Lemma~\ref{lem:rss-deriv},
Proposition~\ref{prop:grad}, and Lemma~\ref{lem:kband-deriv}.
In practice, damping or exponential moving averages in \((\beta,\xi)\) and simple box constraints on the parameter space \(\Theta\) yield stable online updates; see the separate Algorithms section for concrete update rules.
\end{remark}

% \subsection{Empirical Bayes estimator for length and dependence}
% \label{sec:eb-theory}

% % --- Safe theorem/proof envs if not already defined in the preamble
% \makeatletter
% \@ifundefined{theorem}{\newtheorem{theorem}{Theorem}[section]}{}
% \@ifundefined{lemma}{\newtheorem{lemma}[theorem]{Lemma}}{}
% \@ifundefined{proposition}{\newtheorem{proposition}[theorem]{Proposition}}{}
% \@ifundefined{corollary}{\newtheorem{corollary}[theorem]{Corollary}}{}
% \@ifundefined{assumption}{\newtheorem{assumption}[theorem]{Assumption}}{}
% \@ifundefined{definition}{\newtheorem{definition}[theorem]{Definition}}{}
% \@ifundefined{remark}{\newtheorem{remark}[theorem]{Remark}}{}
% \makeatother
% \ifx\proof\undefined
%   \newenvironment{proof}{\par\noindent\textit{Proof.}\ }{\hfill$\square$\par}
% \fi

% We consider trajectories \(i=1,\dots,G\) with length \(L_i\in\mathbb N\) and scalarized return \(g_i\in\mathbb R\).
% Adopt the Normal working model
% \begin{equation}
% \label{eq:model}
% g_i \mid \mu,\sigma^2,\vartheta \sim \mathcal N\!\big(\mu,\; v_i(\vartheta)\big),
% \qquad
% v_i(\vartheta)=\sigma^2\,L_i^{\beta}\,s(L_i;\xi),
% \end{equation}
% where \(\mu\in\mathbb R\) is the common location, \(\sigma^2>0\) is a global scale,
% \(\beta\in\mathbb R\) (the length exponent) controls the length–variance law,
% and \(s(L;\xi)>0\) captures within-trajectory dependence via a low-dimensional shape \(\xi\).
% Define the precisions and normalized weights
% \begin{equation}
% \label{eq:prec-weights}
% \omega_i(\beta,\xi)=\frac{1}{L_i^{\beta}\,s(L_i;\xi)},
% \qquad
% w_i(\beta,\xi)=\frac{\omega_i(\beta,\xi)}{\sum_{j=1}^G \omega_j(\beta,\xi)}.
% \end{equation}
% We take \(p(\mu)\propto 1\) and Jeffreys \(p(\sigma^2)\propto 1/\sigma^2\); weak priors \(\pi_{\beta},\pi_{\xi}\) may be used or set flat for EB/type-II ML.

% \paragraph{Weighted aggregator.}
% The Bayes estimator of \(\mu\) under \eqref{eq:model} is the precision-weighted mean
% \begin{equation}
% \label{eq:mu-agg}
% m(\beta,\xi)=\sum_{i=1}^G w_i(\beta,\xi)\,g_i
% =\frac{\sum_{i=1}^G \omega_i(\beta,\xi)\,g_i}{\sum_{i=1}^G \omega_i(\beta,\xi)}.
% \end{equation}
% Let \(e_i(\beta,\xi)=g_i-m(\beta,\xi)\),
% \(\Lambda_\omega(\beta,\xi)=\sum_{i=1}^G \omega_i(\beta,\xi)\), and
% \begin{equation}
% \label{eq:Lambda-RSS}
% \mathrm{RSS}_\omega(\beta,\xi)=\sum_{i=1}^G \omega_i(\beta,\xi)\,e_i(\beta,\xi)^2.
% \end{equation}

% \begin{proposition}[Posterior in closed form, conditional on \((\beta,\xi)\)]
% \label{prop:posterior}
% Under \eqref{eq:model} with \(p(\mu)\propto 1\) and \(p(\sigma^2)\propto \frac{1}{\sigma^2}\),
% \[
% \mu \mid \sigma^2,\mathbf g,\vartheta \sim \mathcal N\!\big(m(\beta,\xi),\,\sigma^2/\Lambda_\omega(\beta,\xi)\big),
% \quad
% \sigma^2 \mid \mathbf g,\vartheta \sim \mathrm{Inv\!-\!Gamma}\!\left(\frac{G-1}{2},\,\frac{\mathrm{RSS}_\omega(\beta,\xi)}{2}\right).
% \]
% \end{proposition}
% \begin{proof}
% Complete the square in \(\mu\) and integrate the Normal kernel to obtain the Normal posterior for \(\mu\) with variance \(\sigma^2/\Lambda_\omega\).
% Collecting \(\sigma^2\)-powers and exponentials yields an inverse-gamma kernel with shape \((G-1)/2\) and scale \(\mathrm{RSS}_\omega/2\).
% \end{proof}

% \begin{proposition}[GLS optimality of the aggregator]
% \label{prop:gls}
% Fix \((\beta,\xi)\). Among all linear unbiased estimators \(\tilde\mu=\sum_i a_i g_i\) with \(\sum_i a_i=1\),
% \(m(\beta,\xi)\) in \eqref{eq:mu-agg} attains the minimum variance \(\mathrm{Var}(\tilde\mu)=\sigma^2/\Lambda_\omega(\beta,\xi)\).
% \end{proposition}
% \begin{proof}
% This is the Gauss–Markov/GLS result with known heteroscedastic variances
% \(\mathrm{Var}(g_i\mid\mu,\vartheta)=\sigma^2/\omega_i\): the BLUE uses weights proportional to \(\omega_i\), giving variance \(\sigma^2/\sum_i \omega_i\).
% \end{proof}

% \begin{theorem}[EB marginal objective for \((\beta,\xi)\)]
% \label{thm:EB-obj}
% Under \eqref{eq:model} with the above priors, the marginal type-II log-posterior is (up to an additive constant)
% \begin{equation}
% \label{eq:EB-obj}
% \ell(\beta,\xi)
% =\log \pi_\beta(\beta)+\log \pi_\xi(\xi)
% +\tfrac{1}{2}\sum_{i=1}^G \log \omega_i(\beta,\xi)
% -\tfrac{1}{2}\log \Lambda_\omega(\beta,\xi)
% -\tfrac{G-1}{2}\log \mathrm{RSS}_\omega(\beta,\xi).
% \end{equation}
% \end{theorem}
% \begin{proof}
% Integrate \(\mu\) using completion of the square:
% \(\sum_i \omega_i(g_i-\mu)^2=\Lambda_\omega(\mu-m)^2+\mathrm{RSS}_\omega\).
% The \(\mu\)-integral contributes \(-\tfrac{1}{2}\log\Lambda_\omega\).
% Integrate \(\sigma^2\) with Jeffreys prior to obtain a factor proportional to \(\mathrm{RSS}_\omega^{-(G-1)/2}\).
% Collect \(\prod_i \omega_i^{1/2}\) from the Normalizer. Take logs and add priors.
% \end{proof}

% \begin{lemma}[Weighted RSS derivative identity]
% \label{lem:rss-deriv}
% For any component \(\phi\in\{\beta\}\cup\{\xi\text{ components}\}\),
% \begin{equation}
% \label{eq:rss-deriv}
% \frac{\partial}{\partial \phi}\,\mathrm{RSS}_\omega(\beta,\xi)
% =\sum_{i=1}^G \frac{\partial \omega_i}{\partial \phi}\,e_i^2
% =\sum_{i=1}^G \omega_i\,e_i^2\,\frac{\partial \log \omega_i}{\partial \phi}.
% \end{equation}
% \end{lemma}
% \begin{proof}
% Differentiate \(\sum_i \omega_i (g_i-m)^2\). The cross-term \(2\sum_i \omega_i (g_i-m)(-\partial m)\) vanishes because \(\sum_i \omega_i (g_i-m)=0\) by the definition of \(m\).
% \end{proof}

% \begin{proposition}[Gradient of the EB objective]
% \label{prop:grad}
% With \(\Lambda_\omega,\mathrm{RSS}_\omega\) as above,
% \begin{equation}
% \label{eq:EB-grad-general}
% \frac{\partial \ell}{\partial \phi}
% =\frac{1}{2}\!\left[
% \sum_{i=1}^G \frac{\partial \log \omega_i}{\partial \phi}
% -\frac{1}{\Lambda_\omega}\sum_{i=1}^G \frac{\partial \omega_i}{\partial \phi}
% \right]
% -\frac{G-1}{2}\frac{1}{\mathrm{RSS}_\omega}
% \sum_{i=1}^G \omega_i e_i^2\,\frac{\partial \log \omega_i}{\partial \phi}
% +\frac{\partial \log \pi(\phi)}{\partial \phi}.
% \end{equation}
% \end{proposition}
% \begin{proof}
% Differentiate \eqref{eq:EB-obj}; use \(\partial\log\Lambda_\omega/\partial\phi=\Lambda_\omega^{-1}\sum_i \partial\omega_i/\partial\phi\) and Lemma~\ref{lem:rss-deriv}.
% \end{proof}

% \begin{corollary}[Gradient in the length direction]
% \label{cor:grad-beta}
% Since \(\log\omega_i=-\beta\log L_i-\log s(L_i;\xi)\),
% \(\partial \log\omega_i/\partial\beta=-\log L_i\) and \(\partial \omega_i/\partial\beta=-\omega_i \log L_i\).
% Hence
% \begin{equation}
% \label{eq:EB-grad-beta}
% \frac{\partial \ell}{\partial \beta}
% =\frac{1}{2}\!\left[-\sum_{i=1}^G \log L_i
% +\frac{\sum_{i=1}^G \omega_i \log L_i}{\Lambda_\omega}\right]
% +\frac{G-1}{2}\frac{1}{\mathrm{RSS}_\omega}\sum_{i=1}^G \omega_i e_i^2 \log L_i
% +\frac{\partial \log \pi_\beta(\beta)}{\partial \beta}.
% \end{equation}
% \end{corollary}

% \begin{definition}[Stretched-geometric \(k\)-band dependence]
% \label{def:kband}
% Given \(k\in\mathbb N\), \(\rho\in(-1,1)\), \(\eta\ge 0\), define with \(m=\min\{k,L-1\}\)
% \begin{equation}
% \label{eq:s-family}
% s(L;\rho,k,\eta)=1+\frac{2}{L}\sum_{h=1}^m (L-h)\,\rho^{\,h^{\eta}}.
% \end{equation}
% \end{definition}

% \begin{lemma}[Derivatives of the \(k\)-band shape]
% \label{lem:kband-deriv}
% Let \(s_i=s(L_i;\rho,k,\eta)\), \(m_i=\min\{k,L_i-1\}\). Then
% \[
% \frac{\partial \log \omega_i}{\partial \rho}
% =-\frac{1}{s_i}\frac{\partial s_i}{\partial \rho},
% \quad
% \frac{\partial s_i}{\partial \rho}
% =\frac{2}{L_i}\sum_{h=1}^{m_i} (L_i-h)\, h^{\eta}\, \rho^{\,h^{\eta}-1},
% \]
% \[
% \frac{\partial \log \omega_i}{\partial \eta}
% =-\frac{1}{s_i}\frac{\partial s_i}{\partial \eta},
% \quad
% \frac{\partial s_i}{\partial \eta}
% =\frac{2}{L_i}\sum_{h=1}^{m_i} (L_i-h)\, \rho^{\,h^{\eta}} (\log \rho)\, h^{\eta}\log h.
% \]
% \end{lemma}
% \begin{proof}
% Differentiate \eqref{eq:s-family} termwise; use \(\partial\log\omega_i/\partial\cdot=-\partial\log s_i/\partial\cdot\).
% \end{proof}

% \begin{proposition}[Unbiasedness of \(m(\beta,\xi)\)]
% \label{prop:unbiased}
% For fixed \((\beta,\xi)\), \(m(\beta,\xi)\) in \eqref{eq:mu-agg} is unbiased for \(\mu\).
% \end{proposition}
% \begin{proof}
% \(\mathbb E[g_i\mid\mu,\vartheta]=\mu\). Since \(\sum_i w_i=1\), \(\mathbb E[m]=\sum_i w_i\mu=\mu\).
% \end{proof}

% \begin{proposition}[Link to the \(\Delta L\) rule]
% \label{prop:deltaL}
% If \(s(L;\xi)\equiv 1\) and \(\beta=1\), then \(\omega_i=1/L_i\) and \(w_i \propto 1/L_i\), recovering the \(\Delta L\) baseline weighting.
% \end{proposition}
% \begin{proof}
% Immediate from \eqref{eq:prec-weights}.
% \end{proof}

% \begin{assumption}[Regularity and identifiability]
% \label{assump:reg}
% (i) \(\{(g_i,L_i)\}\) are i.i.d.\ or strictly stationary ergodic with \(\sup_i \mathbb E[g_i^4]<\infty\).
% (ii) \(\Theta=\{\beta\}\times\Xi\) is compact (\(\beta\in[\beta_{\min},\beta_{\max}]\), \(|\rho|\le \rho_{\max}<1\), \(\eta\in[0,\eta_{\max}]\), fixed \(k\)).
% (iii) \(\ell_\infty(\vartheta)=\mathbb E[\ell(\vartheta;\mathbf g)]\) has a unique maximizer \(\vartheta^\star\).
% (iv) \(s(L;\xi)\) is smooth in \(\xi\) and bounded away from \(0\) and \(\infty\) on \(\mathrm{supp}(L)\times\Xi\).
% \end{assumption}

% \begin{theorem}[Consistency \& asymptotic normality of the EB estimator]
% \label{thm:consistency}
% Let \(\widehat\vartheta_G=\arg\max_{\vartheta\in\Theta}\ell_G(\vartheta)\).
% Under Assumption~\ref{assump:reg}, \(\widehat\vartheta_G \xrightarrow{p} \vartheta^\star\) as \(G\to\infty\).
% If additionally \(\ell_\infty\) is twice continuously differentiable at \(\vartheta^\star\) with negative-definite Hessian \(H=-\nabla^2\ell_\infty(\vartheta^\star)\) and \(\mathbb E[\|\nabla \ell_1(\vartheta^\star)\|^2]<\infty\),
% then
% \[
% \sqrt{G}(\widehat\vartheta_G-\vartheta^\star)\ \xrightarrow{d}\ \mathcal N\!\big(0,\; H^{-1}\,J\,H^{-1}\big),\qquad
% J=\mathrm{Var}\!\big(\nabla \ell_1(\vartheta^\star)\big).
% \]
% \end{theorem}
% \begin{proof}
% Uniform LLN on compact \(\Theta\) with the stated moments gives \(\ell_G\to\ell_\infty\) uniformly; uniqueness of the maximizer yields argmax consistency (standard M-estimation).
% A Taylor expansion of \(\nabla \ell_G\) at \(\vartheta^\star\) together with a CLT for \(\nabla \ell_1(\vartheta^\star)\) yields the sandwich limit.
% \end{proof}

% \begin{remark}[Smooth but not globally concave]
% While \(\ell(\beta,\xi)\) is not globally concave, \(\nabla \ell\) and the local Hessian are explicit via Lemma~\ref{lem:rss-deriv}, Proposition~\ref{prop:grad}, and Lemma~\ref{lem:kband-deriv}.
% In practice, damping/EMA and box constraints produce stable online updates; see the separate Algorithms section.
% \end{remark}

% \subsection{Weights under $k$-banded dependence with parameter $\rho$}\label{subsec:kband-rho}

% We work with the scalarized series $Y_{i,t}:=\bm{u}^\top \boldsymbol{Z}_{i,t}$ for a fixed unit vector $\bm{u}$ as in \eqref{eq:def_gi_scalar}, so that
% \[
% v_i \;=\; \mathrm{Var}(g_i)\;=\; \sum_{h=-(L_i-1)}^{L_i-1}(L_i-|h|)\,\gamma(h),
% \qquad \gamma(h):=\mathrm{Cov}(Y_{i,t},Y_{i,t-h}).
% \]
% Assume second–order stationarity and the \emph{$k$-banded, geometric} dependence template
% \begin{equation}
% \gamma(0)=\tau^2,\qquad
% \gamma(h)=\tau^2\,\rho^{|h|}\ \mathbf{1}\{|h|\le k\},
% \qquad |\rho|<1,
% \label{eq:kband-geom}
% \end{equation}
% so correlations decay geometrically up to lag $k$ and vanish thereafter.

% \begin{lemma}[Dependence shape factor for geometric $k$-band]\label{lem:sL-kband}
% Let $m_i:=\min\{k,L_i-1\}$. Under \eqref{eq:kband-geom},
% \begin{equation}
% v_i \;=\; \tau^2\!\left[
%   L_i \;+\; 2\sum_{h=1}^{m_i} (L_i-h)\,\rho^{h}
% \right]
% \;=\; \tau^2\,L_i\,s(L_i;\rho,k),
% \qquad
% s(L;\rho,k)\;=\;1+\frac{2}{L}\sum_{h=1}^{m} (L-h)\rho^{h}.
% \label{eq:v-and-s-kband}
% \end{equation}
% Moreover, the sum admits the closed form
% \begin{equation}
% s(L;\rho,k)
% \;=\;
% 1+\frac{2}{L}\Bigg(
%   L\underbrace{\frac{\rho\,(1-\rho^{m})}{1-\rho}}_{S_1(m)}
%   \;-\;
%   \underbrace{\frac{\rho\big(1-(m+1)\rho^{m}+m\rho^{m+1}\big)}{(1-\rho)^2}}_{S_2(m)}
% \Bigg),
% \qquad m=\min\{k,L-1\}.
% \label{eq:s-closed-form}
% \end{equation}
% \end{lemma}

% \begin{proof}
% By Bartlett’s identity (Lemma~\ref{lem:bartlett} in your text),
% \(
% v_i=\sum_{h=-(L_i-1)}^{L_i-1}(L_i-|h|)\gamma(h)
% \).
% Insert \eqref{eq:kband-geom}, collect symmetric lags, and restrict to $|h|\le m_i$:
% \[
% v_i=\tau^2\left[L_i + 2\sum_{h=1}^{m_i}(L_i-h)\rho^h\right].
% \]
% Define $s(L_i;\rho,k):=v_i/(\tau^2 L_i)$ to obtain \eqref{eq:v-and-s-kband}. The closed form \eqref{eq:s-closed-form} follows from
% \(
% \sum_{h=1}^{m}\rho^h=\rho(1-\rho^m)/(1-\rho)
% \)
% and
% \(
% \sum_{h=1}^{m} h\rho^h=\rho\big(1-(m+1)\rho^{m}+m\rho^{m+1}\big)/(1-\rho)^2
% \).
% \end{proof}

% \paragraph{Asymptotics (fixed $k$).}
% For $L\to\infty$ with fixed $k$,
% \begin{equation}
% s(L;\rho,k)
% \;=\;
% 1+2\sum_{h=1}^{k}\rho^h \;-\; \frac{2}{L}\sum_{h=1}^{k} h\rho^h
% \;=\;
% \underbrace{1+\frac{2\rho(1-\rho^{k})}{1-\rho}}_{=:s_\infty(\rho,k)}
% \;+\; O\!\left(\frac{1}{L}\right),
% \label{eq:s-asymp}
% \end{equation}
% so the dependence correction tends to a constant $s_\infty(\rho,k)>0$ (for $|\rho|<1$).

% \begin{proposition}[Weighting rule under geometric $k$-band]\label{prop:weights-kband}
% Under the heteroscedastic model \eqref{eq:hetero_general_prior} with variance law
% \(
% v_i(\vartheta)=\sigma^2 L_i^{\beta}\,s(L_i;\rho,k)
% \),
% the posterior-precision (Bayes) weights satisfy, up to a common factor,
% \begin{equation}
% x_i^{\mathrm{Bayes}} \ \propto\ \frac{1}{L_i^{\beta}\,s(L_i;\rho,k)}.
% \label{eq:weights-kband}
% \end{equation}
% In particular, under the Jeffreys baseline and length exponent $\beta=1$,
% \begin{equation}
% x_i \ \propto\ \frac{1}{\,L_i\,\Big[\,1+\frac{2}{L_i}\sum_{h=1}^{m_i}(L_i-h)\rho^{h}\,\Big]}
% \ =\ \frac{1}{\,L_i + 2\sum_{h=1}^{m_i}(L_i-h)\rho^{h}\,}\!,
% \qquad m_i=\min\{k,L_i-1\}.
% \label{eq:weights-kband-beta1}
% \end{equation}
% For $L_i\to\infty$ with fixed $k$, this is asymptotically
% \(
% x_i \propto \big[L_i\,s_\infty(\rho,k)\big]^{-1}
% \),
% so the length exponent is unchanged and dependence acts as a multiplicative constant.
% \end{proposition}

% \begin{proof}
% From \eqref{eq:hetero_general_prior} with the $k$-banded shape $s(\cdot)$, precision factorizes as
% \(
% \lambda_i(\vartheta)=1/v_i(\vartheta)=(1/\sigma^2)\,L_i^{-\beta}\,s(L_i;\rho,k)^{-1}
% \).
% By Theorem~\ref{thm:ppa}, Bayes weights are proportional to the posterior mean precision; the nuisance scale $1/\sigma^2$ cancels upon normalization (Proposition: Scale cancellation), yielding \eqref{eq:weights-kband}. The $\beta=1$ specialization gives \eqref{eq:weights-kband-beta1}. The asymptotic statement follows from \eqref{eq:s-asymp}.
% \end{proof}

% \paragraph{Variant: flat correlation within the $k$-band.}
% If instead $\gamma(h)=\tau^2\rho\,\mathbf{1}\{1\le |h|\le k\}$ (same correlation for all lags up to $k$), then with $m=\min\{k,L-1\}$,
% \[
% s(L;\rho,k) \;=\; 1 + \frac{2}{L}\sum_{h=1}^{m}(L-h)\rho
% \;=\; 1 + 2m\rho - \frac{m(m+1)}{L}\rho,
% \]
% and the weights are again $x_i \propto \big[L_i^{\beta}\,s(L_i;\rho,k)\big]^{-1}$. (Positive definiteness places constraints on $\rho$ in this template; the geometric form \eqref{eq:kband-geom} with $|\rho|<1$ is always safe.)


\subsection{Weights under non-vanishing constant-lag covariance (compound symmetry)}\label{subsec:cs}

We work with the scalarized token series $Y_{i,t}:=\bm{u}^\top\boldsymbol{Z}_{i,t}$ as in \eqref{eq:def_gi_scalar}. Assume second–order stationarity and
\begin{equation}
  \gamma(0)=\mathrm{Var}(Y_{i,t})=\tau^2,\qquad
  \gamma(h)=\mathrm{Cov}(Y_{i,t},Y_{i,t-h})=\tau^2\rho\quad\text{for all } h\neq 0,
  \label{eq:cs-cov}
\end{equation}
with $\rho\in(-1/(L_i-1),\,1)$ to ensure positive definiteness for a given length $L_i$. (For uniformly valid models across unbounded $L$, take $\rho\in[0,1)$.)

\begin{lemma}[Trajectory variance under compound symmetry]\label{lem:cs-var}
Let $g_i=\sum_{t=1}^{L_i} Y_{i,t}$. Under \eqref{eq:cs-cov},
\begin{equation}
  \mathrm{Var}(g_i)
  \;=\; \tau^2\!\left[\,L_i \;+\; 2\rho\!\sum_{h=1}^{L_i-1} (L_i-h)\right]
  \;=\; \tau^2\,L_i\Big(1+\rho(L_i-1)\Big)
  \;=\; \tau^2\Big((1-\rho)L_i + \rho L_i^2\Big).
  \label{eq:cs-var}
\end{equation}
Equivalently, with the shape factor $s(L;\rho):=\mathrm{Var}(g_i)/(\tau^2 L)$,
\begin{equation}
  s(L;\rho) \;=\; 1+\rho(L-1) \;=\; (1-\rho)+\rho L.
  \label{eq:cs-shape}
\end{equation}
\end{lemma}

\begin{proof}
By Bartlett’s identity (Lemma~\ref{lem:bartlett}),
\[
\mathrm{Var}(g_i)
=\sum_{h=-(L_i-1)}^{L_i-1}(L_i-|h|)\gamma(h)
= L_i\gamma(0)+2\sum_{h=1}^{L_i-1}(L_i-h)\gamma(h).
\]
Insert \eqref{eq:cs-cov} to obtain the first equality in \eqref{eq:cs-var}; the second follows from
$\sum_{h=1}^{L-1}(L-h)=L(L-1)/2$. The shape factor is \eqref{eq:cs-shape}.
\smallskip

(Alternative matrix proof.) The $L_i\times L_i$ covariance matrix of $(Y_{i,1},\dots,Y_{i,L_i})$ is
$\tau^2\!\big((1-\rho)\boldsymbol{I}+\rho\,\boldsymbol{1}\boldsymbol{1}^\top\big)$. Then
$\mathrm{Var}(g_i)=\boldsymbol{1}^\top \mathrm{Cov}\,\boldsymbol{1}
=\tau^2\big((1-\rho)L_i+\rho L_i^2\big)$.
\end{proof}

\begin{proposition}[Weighting rule under constant-lag covariance]\label{prop:cs-weights}
Under the heteroscedastic model \eqref{eq:hetero_general_prior} with
\(
v_i(\vartheta)=\sigma^2 L_i^{\beta}\,s(L_i;\rho)
\)
and $s(L;\rho)$ from \eqref{eq:cs-shape}, the posterior-precision (Bayes) weights satisfy, up to a common factor,
\begin{equation}
  x_i^{\mathrm{Bayes}}
  \ \propto\
  \frac{1}{L_i^{\beta}\,s(L_i;\rho)}
  \;=\;
  \frac{1}{L_i^{\beta}\big(1+\rho(L_i-1)\big)}.
  \label{eq:cs-weights}
\end{equation}
In particular, under the Jeffreys baseline with $\beta=1$,
\begin{equation}
  x_i \ \propto\ \frac{1}{L_i\big(1+\rho(L_i-1)\big)}
  \;=\; \frac{1}{(1-\rho)L_i + \rho L_i^2}.
  \label{eq:cs-weights-beta1}
\end{equation}
\end{proposition}

\begin{proof}
From \eqref{eq:hetero_general_prior} and Lemma~\ref{lem:cs-var},
\(
\lambda_i(\vartheta)=1/v_i(\vartheta)=(1/\sigma^2)\,L_i^{-\beta}\,s(L_i;\rho)^{-1}.
\)
By Theorem~\ref{thm:ppa}, $x_i^{\mathrm{Bayes}}\propto \mathbb{E}[\lambda_i(\vartheta)\mid\{g\}]$; the factor $1/\sigma^2$ cancels upon normalization, yielding \eqref{eq:cs-weights}. The specialization $\beta=1$ gives \eqref{eq:cs-weights-beta1}.
\end{proof}

\begin{corollary}[Asymptotics and effective length law]\label{cor:cs-asymp}
If $\rho>0$ is fixed and $L\to\infty$, then $s(L;\rho)\sim \rho L$ and
\[
v_i(\vartheta)\ \sim\ \sigma^2 L_i^{\beta}\,\rho L_i
\quad\Longrightarrow\quad
v_i(\vartheta)=\Theta(L_i^{\beta+1}).
\]
Under the Jeffreys baseline $\beta=1$, $v_i=\Theta(L_i^2)$ and $x_i\propto L_i^{-2}$ asymptotically. Thus persistent positive correlation induces a quadratic penalty in length relative to the i.i.d.\ case.
\end{corollary}

\begin{remark}[Parameter constraints]
For a fixed length $L$, compound symmetry is positive definite for $\rho\in(-1/(L-1),1)$. To model batches with unbounded $L$, take $\rho\in[0,1)$ so that $s(L;\rho)$ remains positive and the covariance is valid for all $L$.
\end{remark}


% =========================================================
% Plug-in induced-variance estimators and named CAPO variants
% =========================================================
\subsection{Variance-law plug-ins and CAPO variants}\label{subsec:capo-plugins}

\paragraph{From covariance to a one-dimensional target.}
The preceding sections derive a common structural conclusion: for a scalarized trajectory-gradient signal $g$ at length $L$, the effect of within-trajectory token covariance enters the statistically efficient weighting rule only through the \emph{induced variance law}
\begin{equation}
\label{eq:induced-variance-law}
v(L)\;:=\;\mathrm{Var}(g\mid L).
\end{equation}
In the working model \eqref{eq:model} (or its GRPO/grouped extension \eqref{eq:model-grouped}), $v(L)$ takes the parametric form
\begin{equation}
\label{eq:vL-parametric}
v(L;\vartheta)=\sigma^2 L^{\beta}s(L;\xi),\qquad \vartheta=(\beta,\xi),\qquad s(L;\xi)>0.
\end{equation}
Crucially, for the group-based update used in GRPO-style training, the CAPO reweighting requires only a \emph{plug-in} estimate $\widehat v(L)$ (up to a common scale), which is converted into within-group precision weights
\begin{equation}
\label{eq:plugin-weights}
\widehat\omega_{p,i}\ \propto\ \widehat v(L_{p,i})^{-1},
\qquad
\widehat w_{p,i}=\frac{\widehat\omega_{p,i}}{\sum_{j\in\mathcal I_p}\widehat\omega_{p,j}}.
\end{equation}
The remainder of the CAPO update (baseline construction, advantage normalization, and policy optimization) is identical across plug-ins (\S\ref{sec:algorithms}).

\paragraph{Named plug-ins.}
We use the term \emph{CAPO variant} to mean ``CAPO with a specific choice of $\widehat v(L)$''. This section defines four submission-ready plug-ins that cover the key dependence regimes analyzed in \S\ref{subsec:token_dependence}--\S\ref{subsec:cs}.

\subsubsection{L-CAPO: independence / length-only plug-in (includes $\Delta L$)}\label{subsec:l-capo-plugin}
Under the i.i.d. token model (\S\ref{subsec:warmup}), Bartlett’s identity reduces to $\mathrm{Var}(g\mid L)=\tau^2 L$, i.e. $v(L)\propto L$. Equivalently, \eqref{eq:vL-parametric} holds with the \emph{identity} dependence shape $s(L;\xi)\equiv 1$ and the correctly specified exponent $\beta=1$.

\begin{proposition}[L-CAPO as the independence plug-in]\label{prop:lcapo-independence}
Assume i.i.d. tokens so that $\gamma(h)=0$ for $h\ne 0$ and $\mathrm{Var}(g\mid L)=\tau^2 L$. Then the plug-in rule \eqref{eq:plugin-weights} with $\widehat v(L)=L$ yields within-group weights $\widehat w_{p,i}\propto L_{p,i}^{-1}$, i.e. the $\Delta L$ normalization ($\alpha=1$) as the correctly specified independence special case.
More generally, \textbf{L-CAPO} denotes the length-only family obtained by fixing $s(L;\xi)\equiv 1$ and estimating only the length exponent $\beta$ (diagonal EB), so that $\widehat v(L)=L^{\widehat\beta}$ and $\widehat w_{p,i}\propto L_{p,i}^{-\widehat\beta}$.
\end{proposition}
\begin{proof}
Under i.i.d. tokens, $v(L)$ is linear and inverse-variance weights are inverse-length (Lemma~\ref{lem:mvu}). Plugging $\widehat v(L)=L$ into \eqref{eq:plugin-weights} yields $\widehat\omega\propto 1/L$ and hence $\widehat w\propto 1/L$ after within-group normalization.
The length-only generalization with $\widehat v(L)=L^{\widehat\beta}$ follows by the same plug-in mapping.
\end{proof}

\subsubsection{LV-CAPO: empirical-Bayes length\,+\,dependence plug-in}\label{subsec:lvcapo-plugin}
\textbf{LV-CAPO} is the fully structured empirical-Bayes instantiation developed in \S\ref{sec:eb-theory}: it fits $(\beta,\xi)$ by marginal likelihood (either the single-group objective \eqref{eq:EB-obj} or the prompt-group objective \eqref{eq:EB-obj-grouped}) under a chosen structured dependence family $s(L;\xi)$ (e.g., the $k$-banded stretched-geometric model of Definition~\ref{def:kband}). The resulting plug-in is
\begin{equation}
\widehat v_{\mathrm{LV}}(L)=L^{\widehat\beta}\,s(L;\widehat\xi),\qquad \widehat\omega_{p,i}^{\mathrm{LV}}\propto \widehat v_{\mathrm{LV}}(L_{p,i})^{-1}.
\end{equation}
The dependence correction $s(L;\widehat\xi)$ makes LV-CAPO sensitive to correlation-induced variance inflation even when $v(L)$ remains asymptotically linear (summable dependence) or becomes superlinear (persistent dependence).

\subsubsection{CAPO-Q: quadratic induced-variance plug-in (compound symmetry / random intercept)}\label{subsec:capo-q-plugin}
Compound symmetry (\S\ref{subsec:cs}) provides a minimal low-dimensional surrogate for persistent within-trajectory correlation: the covariance between any two distinct token positions is approximately constant. Lemma~\ref{lem:cs-var} implies
\begin{equation}
\label{eq:capo-q-law}
\mathrm{Var}(g\mid L)\ \propto\ (1-\rho)L + \rho L^2,
\end{equation}
and under \eqref{eq:vL-parametric} this becomes $v(L;\beta,\rho)\propto L^{\beta}((1-\rho)+\rho L)= (1-\rho)L^{\beta}+\rho L^{\beta+1}$.
This motivates \textbf{CAPO-Q}, a quadratic plug-in that fits a two-term variance law and then applies inverse-variance weighting.

\begin{proposition}[Outcomes-only quadratic plug-in]\label{prop:capoq-fit}
Fix $\beta$ (e.g., $\beta=1$ for the Jeffreys/independence baseline) and consider the regression model
\begin{equation}
\label{eq:capo-q-reg}
\mathbb E\big[e_{p,i}^2\mid L_{p,i}\big]\ \approx\ a\,L_{p,i}^{\beta}+b\,L_{p,i}^{\beta+1},\qquad a\ge 0,\ b\ge 0,
\end{equation}
where $e_{p,i}=g_{p,i}-m_p$ are within-group residuals from \eqref{eq:mu-agg-grouped}. Let $(\widehat a,\widehat b)$ be any nonnegative least-squares fit of \eqref{eq:capo-q-reg} pooled across groups. Then the quadratic plug-in
\begin{equation}
\label{eq:capo-q-plugin-v}
\widehat v_{\mathrm{Q}}(L)=\widehat a\,L^{\beta}+\widehat b\,L^{\beta+1}
\end{equation}
induces precision weights $\widehat\omega_{p,i}^{\mathrm{Q}}\propto \widehat v_{\mathrm{Q}}(L_{p,i})^{-1}$ and within-group weights $\widehat w_{p,i}^{\mathrm{Q}}$ via \eqref{eq:plugin-weights}.
If the data are generated under compound symmetry with parameter $\rho$ (Lemma~\ref{lem:cs-var}) and correctly specified $\beta$, then $(a,b)$ are proportional to $((1-\rho),\rho)$ and CAPO-Q coincides with the oracle inverse-variance rule (up to scale).
\end{proposition}
\begin{proof}
Lemma~\ref{lem:cs-var} gives $\mathrm{Var}(g\mid L)=\tau^2((1-\rho)L+\rho L^2)$. Under \eqref{eq:vL-parametric}, this is $\sigma^2((1-\rho)L^{\beta}+\rho L^{\beta+1})$ for some scale $\sigma^2$. Thus \eqref{eq:capo-q-reg} is correctly specified with $(a,b)\propto ((1-\rho),\rho)$. Since CAPO weights depend on $v(L)$ only through relative scale within a group, the common multiplicative factor cancels in \eqref{eq:plugin-weights}, yielding oracle weights when $\widehat a,\widehat b$ recover the shape.
\end{proof}

\subsubsection{CAPO-HAC: lag-window (Newey--West) induced-variance plug-in}\label{subsec:capo-hac-plugin}
When token dependence is \emph{short-range} (summable autocovariance), the induced variance can be estimated directly via long-run variance estimators. For a scalarized token increment process $Y_{p,i,t}=\bm u^\top \boldsymbol Z_{p,i,t}$, Bartlett’s identity gives
\begin{equation}
\label{eq:bartlett-vL}
v(L)=\mathrm{Var}\Big(\sum_{t=1}^{L} Y_t\Big)=L\gamma(0)+2\sum_{h=1}^{L-1}(L-h)\gamma(h),\qquad \gamma(h)=\mathrm{Cov}(Y_t,Y_{t-h}).
\end{equation}
\textbf{CAPO-HAC} replaces the parametric $s(L;\xi)$ with a lag-window estimate of the long-run variance. Let $K\ge 1$ be a bandwidth and define pooled lag-$h$ autocovariance estimates (centered within each trajectory)
\begin{equation}
\label{eq:gammahat-pooled}
\widehat\gamma(h)=\frac{\sum_{p,i}\sum_{t=h+1}^{L_{p,i}}\big(Y_{p,i,t}-\bar Y_{p,i}\big)\big(Y_{p,i,t-h}-\bar Y_{p,i}\big)}{\sum_{p,i}(L_{p,i}-h)},\qquad \bar Y_{p,i}=\frac{1}{L_{p,i}}\sum_{t=1}^{L_{p,i}}Y_{p,i,t}.
\end{equation}
Then the Bartlett-window (Newey--West) estimator of the induced variance at length $L$ is
\begin{equation}
\label{eq:hac-vhat}
\widehat v_{\mathrm{HAC}}(L)=L\widehat\gamma(0)+2\sum_{h=1}^{\min\{K,L-1\}}(L-h)\Big(1-\frac{h}{K+1}\Big)\widehat\gamma(h),
\end{equation}
and CAPO-HAC uses $\widehat\omega_{p,i}^{\mathrm{HAC}}\propto \widehat v_{\mathrm{HAC}}(L_{p,i})^{-1}$ with within-group normalization \eqref{eq:plugin-weights}. Under standard mixing conditions and $K\to\infty$ sufficiently slowly, \eqref{eq:hac-vhat} is a consistent estimator of the long-run variance and therefore yields consistent inverse-variance reweighting \citep{newey1987simple,andrews1991heteroskedasticity}.

\begin{remark}[Which plug-in when?]
L-CAPO is the independence/length-only surrogate (good when autocovariances are negligible or when a conservative hedge is preferred). LV-CAPO is the fully structured EB surrogate (best when the chosen $s(L;\xi)$ family matches the dependence). CAPO-Q is the minimal persistent-correlation surrogate, designed to capture superlinear $v(L)$ growth with two parameters. CAPO-HAC is a nonparametric short-range surrogate that targets $v(L)$ directly without committing to a parametric $s(L;\xi)$.
\end{remark}
